% !TEX program = xelatex
% ============================================================
%  机器学习 · 讲义  第四讲：逻辑回归
%  海军深蓝 · 经典学术版式
%  编译：xelatex 04_logistic_regression.tex
%        （无 .aux 时连跑三遍；已有中间文件时两遍即可）
%        或直接用 latexmk -xelatex 04_logistic_regression.tex
% ============================================================
\documentclass[11pt,a4paper,fontset=windows,zihao=-4]{ctexart}

% ---------- 版面 ----------
\usepackage{geometry}
\geometry{a4paper,left=2.55cm,right=2.55cm,top=2.55cm,bottom=2.45cm,
  headheight=18pt,headsep=0.52cm,footskip=0.95cm,
  marginparsep=0.40cm,marginparwidth=2.10cm}

% ---------- 数学与字体 ----------
\usepackage{amsmath,amssymb,amsthm,mathtools}
\usepackage{newpxtext,newpxmath}
\usepackage{microtype}

% ---------- 颜色 ----------
\usepackage[table]{xcolor}
\definecolor{navy}{HTML}{0B2545}      % 海军深蓝（主色）
\definecolor{navy2}{HTML}{33587A}     % 次级蓝
\definecolor{steel}{HTML}{4A6478}     % 灰蓝
\definecolor{grey}{HTML}{6B7280}      % 中性灰
\definecolor{mist}{HTML}{F4F6F9}      % 极浅蓝灰（底色）
\definecolor{gold}{HTML}{8A6A3B}      % 暗金（仅用于提示）
\definecolor{goldlight}{HTML}{C6A15B} % 亮金（深蓝底上的金线）

% ---------- 表格与列表 ----------
\usepackage{booktabs}
\usepackage{array}
\usepackage{multirow}
\usepackage{enumitem}
\setlist[itemize]{leftmargin=1.6em,topsep=0.45em,itemsep=0.28em,parsep=0pt}
\setlist[enumerate]{leftmargin=2.0em,topsep=0.45em,itemsep=0.28em,parsep=0pt}

% ---------- 图形 ----------
\usepackage{tikz}
\usetikzlibrary{calc,positioning,arrows.meta,decorations.pathreplacing,
  patterns,backgrounds,fit,shapes.geometric}
\usepackage{pgfplots}
\pgfplotsset{compat=1.18}

% ---------- 彩色框 ----------
\usepackage[most]{tcolorbox}

% ---------- 页眉页脚 ----------
\usepackage{fancyhdr}
\usepackage{titlesec}
\usepackage{caption}
\usepackage{marginnote}
\usepackage{tocloft}

% ---------- 超链接 ----------
\usepackage[colorlinks=true,linkcolor=navy,urlcolor=steel,citecolor=navy,
  bookmarksnumbered=true]{hyperref}
\hypersetup{
  pdftitle={机器学习·讲义（第四讲）：逻辑回归},
  pdfauthor={CyArk},
  pdfsubject={对数几率、极大似然、IRLS、完全分离与概率校准},
  pdfkeywords={逻辑回归, 对数几率, 交叉熵, 牛顿法, IRLS, 完全分离, 概率校准}
}

% ============================================================
%  数学宏
% ============================================================
\newcommand{\R}{\mathbb{R}}
\newcommand{\bX}{\mathbf{X}}
\newcommand{\by}{\mathbf{y}}
\newcommand{\bw}{\mathbf{w}}
\newcommand{\bx}{\mathbf{x}}
\newcommand{\bz}{\mathbf{z}}
\newcommand{\bg}{\mathbf{g}}
\newcommand{\bp}{\mathbf{p}}
\newcommand{\bv}{\mathbf{v}}
\newcommand{\bu}{\mathbf{u}}
\newcommand{\bs}{\mathbf{s}}
\newcommand{\bA}{\mathbf{A}}
\newcommand{\bB}{\mathbf{B}}
\newcommand{\bD}{\mathbf{D}}
\newcommand{\bI}{\mathbf{I}}
\newcommand{\bH}{\mathbf{H}}
\newcommand{\bS}{\mathbf{S}}
\newcommand{\bV}{\mathbf{V}}
\newcommand{\bW}{\mathbf{W}}
\newcommand{\bZ}{\mathbf{Z}}
\newcommand{\bSigma}{\boldsymbol{\Sigma}}
\newcommand{\bmu}{\boldsymbol{\mu}}
\newcommand{\bDelta}{\boldsymbol{\Delta}}
\newcommand{\btheta}{\boldsymbol{\theta}}
\newcommand{\bzero}{\mathbf{0}}
\newcommand{\bmone}{\mathbf{1}}
\newcommand{\tp}{^{\top}}
\newcommand{\norm}[1]{\lVert #1\rVert}
\newcommand{\abs}[1]{\lvert #1\rvert}
\DeclareMathOperator*{\argmin}{arg\,min}
\DeclareMathOperator*{\argmax}{arg\,max}
\DeclareMathOperator{\tr}{tr}
\DeclareMathOperator{\rank}{rank}
\DeclareMathOperator{\diag}{diag}
\DeclareMathOperator{\Span}{Span}
\DeclareMathOperator{\sign}{sign}
\DeclareMathOperator{\Var}{Var}
\DeclareMathOperator{\Cov}{Cov}
\DeclareMathOperator{\softplus}{softplus}
\newcommand{\dd}{\mathrm{d}}

% ============================================================
%  版式细节
% ============================================================
\linespread{1.30}
\setlength{\parskip}{0.24em}
\setlength{\parindent}{2em}

% 章节样式
\ctexset{
  section/name={},
  section/number=\arabic{section},
}
\titleformat{\section}[block]
  {\normalfont\sffamily\Large\bfseries\color{navy}}
  {\colorbox{navy}{\makebox[1.55em][c]{\color{white}\thesection}}\hspace{0.75em}}{0pt}{}
  [\vspace{0.30em}{\color{navy!30}\titlerule[0.7pt]}]
\titlespacing*{\section}{0pt}{2.4ex plus 1.0ex minus 0.4ex}{1.5ex}

\titleformat{\subsection}[hang]
  {\normalfont\sffamily\large\bfseries\color{navy!88}}
  {\thesubsection}{0.65em}{}
\titlespacing*{\subsection}{0pt}{1.8ex plus 0.6ex}{0.8ex}

\titleformat{\subsubsection}[runin]
  {\normalfont\sffamily\bfseries\color{navy!80}}
  {}{0pt}{}
\titlespacing*{\subsubsection}{0pt}{1.2ex plus 0.4ex}{0.6em}

\titleformat{\paragraph}[runin]
  {\normalfont\sffamily\bfseries\color{navy!88}}
  {}{0pt}{}
\titlespacing*{\paragraph}{0pt}{1.05ex plus 0.35ex}{0.55em}

% 页眉页脚
\pagestyle{fancy}
\fancyhf{}
\renewcommand{\headrulewidth}{0.5pt}
\renewcommand{\headrule}{{\color{navy!40}%
  \hrule height \headrulewidth width \headwidth}}
\fancyhead[L]{\sffamily\footnotesize\color{steel}机器学习\,·\,讲义}
\fancyhead[R]{\sffamily\footnotesize\color{steel}\leftmark}
\fancyfoot[C]{\sffamily\footnotesize\color{navy}\thepage}
\renewcommand{\sectionmark}[1]{\markboth{\thesection\quad #1}{}}
\fancypagestyle{plain}{%
  \fancyhf{}\renewcommand{\headrulewidth}{0pt}%
  \fancyfoot[C]{\sffamily\footnotesize\color{navy}\thepage}}

% 图表标题
\captionsetup{font={small},labelfont={bf,sf,color=navy},labelsep=quad,
  justification=centering,skip=6pt}

% 目录（手工排版标题，避免 \leftmark 捕获颜色命令）
\makeatletter
\renewcommand{\tableofcontents}{%
  \par\begingroup
    \sffamily\Large\bfseries\color{navy}目\quad 录\par
  \endgroup
  \vspace{0.85em}
  \setlength{\parskip}{0pt}%
  \@starttoc{toc}%
}
\makeatother

% 目录条目样式
\renewcommand{\cftsecfont}{\footnotesize\sffamily\bfseries\color{navy}}
\renewcommand{\cftsecpagefont}{\footnotesize\sffamily\bfseries\color{navy}}
\renewcommand{\cftsubsecfont}{\footnotesize\color{navy!85}}
\renewcommand{\cftsubsecpagefont}{\footnotesize\color{navy!85}}
\renewcommand{\cftsubsubsecfont}{\color{steel}}
\renewcommand{\cftsubsubsecpagefont}{\color{steel}}
\renewcommand{\cftsecleader}{\cftdotfill{\cftdotsep}}
\renewcommand{\cftsubsecleader}{\cftdotfill{\cftdotsep}}
\renewcommand{\cftsubsubsecleader}{\cftdotfill{\cftdotsep}}
\setlength{\cftsecnumwidth}{1.9em}
\setlength{\cftsubsecnumwidth}{2.5em}
\setlength{\cftsubsubsecnumwidth}{3.4em}
\setlength{\cftbeforesecskip}{4.0pt}
\setlength{\cftbeforesubsecskip}{0.2pt}
\setlength{\cftbeforesubsubsecskip}{0pt}
\renewcommand{\cftsecaftersnum}{\hspace{0.55em}}
\renewcommand{\cftsubsecaftersnum}{\hspace{0.4em}}
\renewcommand{\cftsubsubsecaftersnum}{\hspace{0.35em}}

% ============================================================
%  提示框样式
% ============================================================
\tcbset{
  basebox/.style={
    enhanced, breakable, sharp corners,
    boxrule=0pt, leftrule=2.2pt, arc=0pt, outer arc=0pt,
    left=0.9em, right=0.9em, top=0.45em, bottom=0.62em,
    before skip=0.85em, after skip=0.85em,
    before upper={\parindent=0pt\parskip=0.35em},
    fonttitle=\sffamily\bfseries\small,
    coltitle=navy, titlerule=0.35pt,
    toptitle=0.32em, bottomtitle=0.5em,
  }
}
\newtcolorbox[auto counter,number within=section]{dfn}[1][]{%
  basebox, colframe=navy, colback=mist, colbacktitle=mist,
  title={定义~\thetcbcounter\quad\textcolor{navy!45}{\textbar}\quad\textcolor{navy}{#1}}}
\newtcolorbox[auto counter,number within=section]{thm}[1][]{%
  basebox, colframe=navy, colback=white, colbacktitle=white,
  title={定理~\thetcbcounter\quad\textcolor{navy!45}{\textbar}\quad\textcolor{navy}{#1}}}
\newtcolorbox[auto counter,number within=section]{prop}[1][]{%
  basebox, colframe=navy2, colback=white, colbacktitle=white,
  title={命题~\thetcbcounter\quad\textcolor{navy!45}{\textbar}\quad\textcolor{navy}{#1}}}
\newtcolorbox[auto counter,number within=section]{exm}[1][]{%
  basebox, colframe=steel, colback=white, colbacktitle=white,
  title={例~\thetcbcounter\quad\textcolor{navy!45}{\textbar}\quad\textcolor{steel}{#1}}}
\newtcolorbox{add}[1][]{%
  basebox, colframe=steel, colback=mist, colbacktitle=mist,
  title={\textcolor{steel}{补充}\quad\textcolor{steel!45}{\textbar}\quad\textcolor{steel}{#1}}}
\newtcolorbox{warn}[1][]{%
  basebox, colframe=gold, colback=gold!4, colbacktitle=gold!4,
  title={\textcolor{gold}{注意}\quad\textcolor{gold!55}{\textbar}\quad\textcolor{gold}{#1}}}
\newtcolorbox{keybox}[1][]{%
  enhanced, breakable, sharp corners, boxrule=0.55pt,
  colframe=navy!55, colback=navy!3, colbacktitle=navy!3,
  left=0.9em, right=0.9em, top=0.45em, bottom=0.62em,
  before skip=1.0em, after skip=1.0em,
  before upper={\parindent=0pt\parskip=0.35em},
  fonttitle=\sffamily\bfseries\small, coltitle=navy,
  titlerule=0.35pt, toptitle=0.3em, bottomtitle=0.5em,
  title={\textcolor{navy}{#1}}}

% 边注
\renewcommand*{\marginfont}{\footnotesize\color{steel}}
\setlength{\marginparpush}{6pt}

% 常用小工具
\newcommand{\hl}[1]{\textcolor{navy}{\textbf{#1}}}
\newcommand{\tabhead}{\sffamily\bfseries\color{navy}}
\tikzset{
  cls1/.style={fill=navy,draw=none},
  cls2/.style={fill=goldlight,draw=none},
  frame/.style={draw=steel!30,line width=0.4pt},
  axisar/.style={->,steel!55,line width=0.5pt},
  lab/.style={font=\scriptsize,text=steel,inner sep=2pt},
}

\begin{document}

% ============================================================
%  封面 · 目录 · 导读
% ============================================================

\begin{titlepage}
\thispagestyle{empty}
\raggedright
% ---- 顶部深蓝色块 + 金色分隔线 + 例题的真实解 ----
\begin{tikzpicture}[remember picture, overlay]
  \fill[navy] (current page.north west) rectangle ([yshift=-12.30cm]current page.north east);
  \fill[goldlight] ([yshift=-12.30cm]current page.north west)
        rectangle ([yshift=-12.37cm]current page.north east);
  % 右侧：6 个样本 + 三条等概率线（本讲 §4 例题的真实解）
  \begin{scope}[shift={([xshift=12.30cm,yshift=-10.28cm]current page.north west)}]
    \draw[white,opacity=0.09,line width=0.35pt] (0,0) rectangle (7.02,4.60);
    % p = 0.1 / 0.9 两条等概率线（细、虚）
    \draw[white,opacity=0.16,line width=0.35pt,densely dashed] (2.774,4.220) -- (6.531,3.324);
    \draw[white,opacity=0.16,line width=0.35pt,densely dashed] (0.357,2.124) -- (6.531,0.652);
    % p = 0.5 决策边界（金线，实）
    \draw[goldlight,line width=0.85pt] (0.357,3.452) -- (6.531,1.980);
    % 类 1：实心白圆
    \foreach \a/\b in {0.62/0.62, 1.93/0.94, 5.09/2.86}{%
      \fill[white,opacity=0.92] (\a,\b) circle (1.55pt);}
    % 类 0：空心金方
    \foreach \a/\b in {4.82/2.38, 6.40/3.18, 3.51/3.98}{%
      \fill[white,opacity=0.08] ($(\a,\b)+(-1.55pt,-1.55pt)$) rectangle ($(\a,\b)+(1.55pt,1.55pt)$);
      \draw[goldlight,opacity=0.95,line width=0.55pt]
        ($(\a,\b)+(-1.55pt,-1.55pt)$) rectangle ($(\a,\b)+(1.55pt,1.55pt)$);}
    % 两个特殊点：$p=0.295$ 与 $p=0.505$
    \draw[white,opacity=0.55,line width=0.4pt] (5.09,2.86) circle (3.4pt);
    \draw[white,opacity=0.55,line width=0.4pt] (4.82,2.38) circle (3.4pt);
    % 线上标签
    \node[font=\tiny,text=white,opacity=0.45,anchor=west] at (6.56,3.30) {$0.1$};
    \node[font=\tiny,text=white,opacity=0.55,anchor=west] at (6.56,1.96) {$0.5$};
    \node[font=\tiny,text=white,opacity=0.45,anchor=west] at (6.56,0.62) {$0.9$};
  \end{scope}
\end{tikzpicture}%
\vspace*{0.90cm}
{\color{white!58!navy}\sffamily\large 机器学习\,\textperiodcentered\,讲义系列\par}
\vspace{0.90cm}
{\color{white}\sffamily\bfseries\fontsize{44}{48}\selectfont 逻辑回归\par}
\vspace{0.62cm}
{\color{white!82!navy}\sffamily\large 概率输出的判别式路线\par}
\vspace{0.55cm}
{\color{goldlight}\rule{3.0cm}{1.1pt}\par}
\vspace{0.46cm}
{\color{white!58!navy}\itshape\large Logistic Regression\par}

\vspace*{3.10cm}
{\color{goldlight}\sffamily\bfseries\Large 第\ 四\ 讲\par}
\vspace{1.05cm}
{\color{steel}\sffamily\footnotesize 本\ 讲\ 涵\ 盖\par}
\vspace{0.26cm}
{\color{grey}\footnotesize 对数几率\ \ \textperiodcentered\ \ 极大似然\ \ \textperiodcentered\ \ 交叉熵
 \ \ \textperiodcentered\ \ 牛顿法与 IRLS\ \ \textperiodcentered\ \ 完全分离\ \ \textperiodcentered\ \ 概率校准\par}

\vspace{3.55cm}
{\color{navy!30}\rule{3.0cm}{0.4pt}\par}
\vspace{0.32cm}
{\color{grey}\footnotesize https://github.com/CyArk-fs/MLnote\par}
\end{titlepage}

% ============================================================
\pagenumbering{roman}
\setcounter{page}{1}

\addtocontents{toc}{\protect\footnotesize}
\tableofcontents
\thispagestyle{plain}
\clearpage

% ============================================================
\pagenumbering{arabic}
\setcounter{page}{1}

\section*{导\quad 读}
\addcontentsline{toc}{section}{导读}
\markboth{导读}{}

上一讲从一个线性打分里读出了一个类别，代价是先假定每一类都是高斯、且共享同一个协方差。这一讲要把那套假定全部丢掉，只保留结论里真正被用到的一条，并回答一个问题：\textbf{同一个打分，凭什么能读成一个概率？}

听上去像是加了一层装饰——类别都判出来了，还要概率做什么？恰恰相反。概率一旦有了，\emph{决策}这件事才真正开始：代价不等时该往哪边偏，先验不同时阈值该挪到哪里，一个人和一台机器给出的 $0.7$ 到底差在哪。而这些问题的答案，都会反过来改变我们对「模型好不好」的判断——一个只输出类别的分类器，你没法问它「有多确定」。

逻辑回归（logistic regression）是这条路线最经典、也最容易被误解的入口。「回归」两个字常让人以为它是线性回归的亲戚，其实两者的共同点只在骨架：都在算 $\bw\tp\bx$（截距已并入 $\bw$）。区别在于这个分数后面接什么——线性回归把它当预测值，逻辑回归把它压进 $(0,1)$ 当概率。整讲的内容，都是围绕这一步展开的：\emph{凭什么用 sigmoid 压}、\emph{参数怎么定}、\emph{定不出来怎么办}、\emph{定出来以后能不能信}。

\begin{add}[它与前两讲的关系]
第二讲的 LDA 与上一讲的高斯判别分析，边界都是超平面，判据却与逻辑回归完全不同：它们先假设每一类长什么样（高斯），再用贝叶斯公式反推；逻辑回归跳过「每一类长什么样」，直接假设后验的形状。所以前两者是\emph{生成式}，逻辑回归是\emph{判别式}。三者的取舍放在 \S7.2，那里会看到：假设对时生成式更省样本，假设错时逻辑回归更抗造——这不是口号，而是两条曲线的交叉。
\end{add}

\subsection*{读完这一讲，你应当能够}

\begin{enumerate}[label=\textcolor{navy}{\bfseries\arabic*.},leftmargin=2.2em]
  \item 说清「硬标签」缺了什么，以及为什么最小化期望风险需要后验概率本身；
  \item 写出几率、对数几率与 sigmoid，并解释为什么 sigmoid 是\emph{反解}出来的、不是挑出来的；
  \item 从伯努利似然推出交叉熵，说明「似然最大」「交叉熵最小」「KL 最小」为什么是同一件事；
  \item 写出梯度与 Hessian，说明为什么这个问题严格凸、局部最优即全局最优；
  \item 把牛顿法重排成 IRLS，指出它每一步都是一次加权最小二乘；
  \item 判断数据在什么情况下\emph{根本不存在}有限解，分清完全分离与拟完全分离，并知道怎么补救；
  \item 解释 $p=0.7$ 什么时候能当「七成把握」用，阈值为什么默认是 $0.5$，以及怎样校准。
\end{enumerate}

\subsection*{预备知识}

\begin{itemize}
  \item \textbf{线性代数}：矩阵乘法、转置、正定、可逆。整讲的矩阵最大不过 $(D+1)\times(D+1)$，例题里只有 $3\times3$。
  \item \textbf{概率统计}：条件概率、贝叶斯公式、伯努利分布、期望。第三讲的多元正态只在本讲 \S2.3 被引用一次，不重推。
  \item \textbf{微积分}：链式法则、一阶与二阶条件、泰勒展开。牛顿法只用到「二次近似」这一条直觉。
\end{itemize}

\subsection*{记号约定}

沿用前两讲的习惯：\textbf{粗体}表示向量或矩阵，\emph{斜体}表示标量。本讲新增的符号集中列在下面，后文不再解释。

\begin{table}[htbp]
\centering
\small
\begin{tabular}{@{}ll@{\hspace{1.6em}}ll@{}}
\toprule
\rowcolor{navy!7}
\tabhead 记号 & \tabhead 含义 & \tabhead 记号 & \tabhead 含义\\
\midrule
$N$ & 样本数 & $D$ & 特征维数\\
$\bx\in\R^{D+1}$ & 特征（末位恒为 $1$） & $\bX\in\R^{N\times(D+1)}$ & 设计矩阵（末列全 $1$）\\
$\bw\in\R^{D+1}$，$b=w^{(D+1)}$ & 权重；$b$ 为截距 & $z_i=\bw\tp\bx_i$ & 线性得分\\
$p_i=\sigma(z_i)$ & 预测概率 & $y_i\in\{0,1\}$ & 标签\\
$s_i=2y_i-1$ & $\pm1$ 标签 & $\ell_i$ & 单样本对数似然\\
$L$ & 平均交叉熵（损失） & $\bg=\nabla L$ & 梯度（平均尺度）\\
$\bH=\nabla^2L$ & Hessian（平均尺度） & $\bDelta=\bH^{-1}\bg$ & 牛顿增量\\
$\bz_{\mathrm w}$ & IRLS 工作响应 & $\bW=\diag(p_i(1-p_i))$ & 权重（$=$ 方差）\\
$\tau$ & 判决阈值（$p$ 尺度） & $z^\ast$ & 判决阈值（$z$ 尺度）\\
$c_{ij}$ & 真类 $i$ 判成 $j$ 的代价 & $\phi_c=p(y=c)$ & 类先验\\
\bottomrule
\end{tabular}
\end{table}

全篇一律用增广记号：$\bx$ 与 $\bw$ 都是 $D+1$ 维，末位分别固定为 $1$ 与截距 $b$；$\bX$ 的末列全为 $1$。于是 $z=\bw\tp\bx$ 里不再单独出现 $b$。需要点名截距时记 $b=w^{(D+1)}$——\S5.3 说「截距不惩罚」、\S6.1 的先验修正都要用到它。

有三处约定必须一次说清，否则后文会处处含混。

\paragraph{第一，尺度。}$\ell_i=y_i\log p_i+(1-y_i)\log(1-p_i)$ 是\textbf{正项}的对数似然，越大越好；
\[
  L=-\frac1N\sum_{i=1}^{N}\ell_i
\]
是\textbf{损失}，也是平均交叉熵，越小越好。梯度与 Hessian 一律取\emph{平均尺度}：
$\bg=\frac1N\sum_i(p_i-y_i)\bx_i$、$\bH=\frac1N\bX\tp\bW\bX$。
牛顿步 $\bw\leftarrow\bw-\bH^{-1}\bg$ 里分子分母同乘 $N$，所以\textbf{步长与是否平均无关}；但\textbf{停止容差受影响}——容差写 $10^{-6}$ 时，平均尺度与合计尺度差 $N$ 倍。这是实现里最常见的「两份代码判据不一致」的来源。

\paragraph{第二，与上一讲的方向差异。}字母沿用第二讲的「投影得分」，含义却反了：
第二讲的 $\hat\bw$ 指向类 2，得分越高越像类 2，判决是 $z>z_0$ 取类 2；
本讲的 $z$ 是后验对数几率，\textbf{越大越像类 1}。
两讲各自自洽，但引用阈值与先验修正时\textbf{只能借结构、不能借符号}。

\paragraph{第三，$\delta$ 与 $\bDelta$ 的分工。}$\delta_c(\bx)$ 沿用第二讲，是判别函数；
本讲的牛顿增量记 $\bDelta$。工作响应记 $\bz_{\mathrm w}$（$N$ 维向量），与单样本得分 $z_i$ 靠下标区分。

\subsection*{四种读法}

\begin{itemize}
  \item \textbf{只想弄懂「为什么要概率、sigmoid 从哪来」}：读 \S1、\S2，再看 \S4.5 的算例。
  \item \textbf{要动手实现}：\S3.3 给梯度、\S4.3 给 IRLS、\S4.4 给数值细节，三节连起来就是一份可用的代码说明。
  \item \textbf{被「算不出来」坑过}：直接看 \S5，那里有反例、有判据、也有三种补救。
  \item \textbf{关心模型可信不可信}：\S6 全读，重点是校准与指标两节。
\end{itemize}

% ============================================================
%  第 1 节
% ============================================================

\section{概率输出的必要性}

\subsection{硬标签的代价}

三个分类器摆在桌上：感知机、LDA、还有本讲的逻辑回归。前两个都只吐出一个类别（或一个符号），差别只在「怎么划这条线」。问题是：\emph{划完线以后呢？}

先看感知机。它的输出是 $\sign(\bw\tp\bx)$，损失是「错几个」。这个设计有个隐含前提：\textbf{所有错误一样严重}。可现实里不是。垃圾邮件把老板的信拦了，和促销邮件漏进收件箱，代价差着数量级；体检模型漏诊一个早期病灶，与多做一个复查，更不是一回事。用一个只数错误个数的模型，你没法表达这件事。

再看 LDA。它比感知机多给了一样东西：一个可以比较的得分 $z=\bw\tp\bx$。既然有得分，能不能把「$z$ 离边界多远」当把握程度？

不能——因为\textbf{这个得分没有单位}。用第二讲 \S2.5 的例题数据：$\hat\bw=(0.622,\,0.783)\tp$（第二讲的记号，二维、不含截距位），十个样本的得分落在 $[3.41,4.44]$ 与 $[5.30,5.40]$ 两段。现在把 $\bw$ 整体乘 $10$，得分全部变成原来的十倍，阈值从 $4.63$ 变成 $46.3$——判决结果一个不变。

\begin{center}
\small
\begin{tabular}{@{}lccc@{}}
\toprule
\rowcolor{navy!7}
\tabhead & \tabhead 类 1 得分区间 & \tabhead 类 2 得分区间 & \tabhead 判决\\
\midrule
$\hat\bw$ & $[3.41,\,4.44]$ & $[5.30,\,5.40]$ & 全部分对\\
$10\hat\bw$ & $[34.1,\,44.4]$ & $[53.0,\,54.0]$ & 全部分对\\
\bottomrule
\end{tabular}
\end{center}

同一份数据、同一个方向，得分翻十倍而结论不变。所以 $z$ 的\emph{绝对值}不携带任何信息，只有它的\emph{符号}（在哪一侧）有用。想让它携带信息，必须给它一个统一的刻度。

\subsection{期望风险与后验}

刻度从哪里来？不从几何来，从代价来。

\begin{dfn}[后验概率与期望风险]
设 $p=p(y=1\mid\bx)$ 是给定特征后样本属于类 1 的后验概率。记 $c_{ij}$ 为\textbf{真类 $i$ 判成 $j$} 的代价，且只罚错判：$c_{00}=c_{11}=0$，$c_{01},c_{10}>0$。两种判决的期望风险分别是
\begin{equation}
  R(\text{判 }1)=c_{01}\,(1-p),\qquad
  R(\text{判 }0)=c_{10}\,p .
  \label{eq:risk}
\end{equation}
使期望风险最小的判决规则称为\textbf{贝叶斯判决}。
\end{dfn}

式 \eqref{eq:risk} 值得逐项读一遍。判 1 时，只有真实为 0 的样本会付出 $c_{01}$，这类样本占 $1-p$；判 0 时，只有真实为 1 的样本付出 $c_{10}$，占 $p$。两行都没有 $c_{00}$ 与 $c_{11}$，因为判对不罚。

\begin{prop}[贝叶斯判决由后验与代价比决定]
在式 \eqref{eq:risk} 的记号下，判 1 当且仅当
\begin{equation}
  p(y=1\mid\bx)\;>\;\tau,\qquad \tau=\frac{c_{01}}{c_{01}+c_{10}} .
  \label{eq:bayesrule}
\end{equation}
等代价（$c_{01}=c_{10}$）时 $\tau=1/2$。
\end{prop}

\begin{proof}
$c_{01}(1-p)<c_{10}p\iff c_{01}<(c_{01}+c_{10})p\iff p>\dfrac{c_{01}}{c_{01}+c_{10}}$。
\end{proof}

这条命题就是「要概率」的硬理由：右边的式子需要 $p$ 本身，而不只是它的符号。只给符号的分类器，无论怎么调阈值，都只能表达「等代价」这一种情形；一旦代价不等，就必须知道后验的\emph{数值}。

这也顺带解释了阈值的来历：$\tau$ 不是 $1/2$ 这个「自然数」，而是代价比算出来的。$c_{01}$ 越大（漏判越贵），$\tau$ 越小，越倾向于判 1。\S6.1 会把这条式子放到先验不等的情形里重讲一遍。

\begin{warn}[「得分大 = 把握大」是错的]
$z$ 大只说明「离边界远」，而「远」是多少厘米取决于 $\norm{\bw}$。同一个样本，把 $\bw$ 乘 $10$ 就能让 $z$ 大十倍，把握却一点没变。要谈把握，先要给 $z$ 一个统一的刻度——这正是 \S2 要做的事。
\end{warn}

\begin{keybox}[本节结论]
\textbf{概率是决策的输入，不是模型的装饰。}
只要代价不等、先验不等，或者你打算把输出交给下游做加权决策，就需要后验概率的数值；而一个只输出符号的分类器，连「阈值该放在哪」这个问题都提不出来。
\end{keybox}

% ============================================================
%  第 2 节
% ============================================================

\section{对数几率与 sigmoid}

\subsection{几率与对数几率}

要给概率一个「统一刻度」，最直接的想法是把 $(0,1)$ 上的 $p$ 拉到整条实轴上。可拉法有无数种，凭什么挑一种？要回答它，得依次说清三件事：为什么从\emph{比值}入手，为什么还要\emph{取一次对数}，为什么不能停在概率或几率上。定义放在最后。

\paragraph{比值的必要性。}先看两个理由。

其一，分类问题的比较天然是\textbf{倍数}，不是差值。「这类比那类可能两倍」在任何基线上都成立；「概率高 $0.2$」只在某个基线上成立，而且会撞墙——$0.5\to0.7$ 说得通，$0.9\to1.1$ 根本不存在。

其二，正负两类必须\textbf{对称}：把两类互换，判断应当恰好反号。差值做不到这一点（$0.9-0.5$ 与 $0.5-0.1$ 都是 $0.4$，方向却相反），比值可以。于是把「发生」与「不发生」的比拿出来，作为第一个候选刻度：
\begin{equation}
  \mathrm{odds}=\frac{p}{1-p}\in(0,+\infty).
\end{equation}
$p=0.5$ 时几率为 $1$，$p=0.9$ 时几率为 $9$（「九比一」，日常语言里的赔率），$p=0.1$ 时几率为 $1/9$。互换两类，几率取倒数——这正是「对称」的样子。

\paragraph{对数的必要性。}取对数才是关键，理由有三层，一层比一层硬。

\textbf{（一）尺度不对称。}几率拉开了 $(0,1)$，却只拉开一半：下界 $0$ 仍是一道墙，上限却跑到 $+\infty$。于是「支持类 1」的证据可以无限大，「支持类 0」的证据最多压到 $0$，两头不等长。取对数后值域是整个 $\R$，$p=1/2$ 落在原点，两侧等距：正负号只说「支持哪一类」，绝对值说「支持多强」。

\textbf{（二）结构不匹配。}这才是真正的原因：几率是\textbf{乘性}的量：两条独立证据合成时，后验几率等于先验几率乘上两个似然比，
\begin{equation}
  \frac{p}{1-p}=\frac{\phi}{1-\phi}\times\mathrm{LR}_1\times\mathrm{LR}_2 .
  \label{eq:oddsmult}
\end{equation}
而我们要用的模型是\textbf{加法}的（$z=\sum_j w^{(j)}x^{(j)}+b$）。乘法塞不进加法的框架，唯一的接法是先取对数：$\ln$ 把式 \eqref{eq:oddsmult} 变成三项相加，于是每个特征恰好对应一项加数，「各条证据的贡献相加」这句话才有了着落。

\textbf{（三）唯一性。}「取对数」也不是在众多函数里选一个。设 $L:(\R_{>0},\times)\to(\R,+)$ 连续且把乘法变成加法，即 $L(ab)=L(a)+L(b)$，则
\begin{equation}
  L(a)=c\ln a,\qquad c=L'(1)\ne0 .
  \label{eq:logunique}
\end{equation}
证明只需三行：对 $b$ 求导得 $aL'(ab)=L'(b)$，令 $b=1$ 得 $L'(a)=L'(1)/a$，再积分并用 $L(1)=0$（由 $L(1)=L(1\cdot1)=2L(1)$）即得。把 $a=\frac{p}{1-p}$ 代进去，得到 $\mathrm{logit}(p)=c\ln\frac{p}{1-p}$；常数 $c$ 只改刻度的疏密，不改排序，也不改判决。\textbf{所以「取对数」不是偏好，而是「要求乘性证据可加」的唯一结果。}

三层理由可以用一个数值例子一次看完。先验 $p=0.5$（几率 $1$），证据 A 的似然比是 $9$，证据 B 是 $4$：

\begin{center}
\small
\begin{tabular}{@{}lll@{}}
\toprule
\rowcolor{navy!7}
\tabhead 尺度 & \tabhead 合成方式 & \tabhead 结果\\
\midrule
概率 $p$ & 没有合成律 & $0.9$ 与 $0.8$ 既不能加也不能乘\\
几率 $\mathrm{odds}$ & 相乘：$1\times9\times4$ & $36$，即 $p=36/37=0.973$\\
对数几率 $z$ & 相加：$0+\ln9+\ln4$ & $3.583$，即 $\sigma(3.583)=0.973$\\
\bottomrule
\end{tabular}
\end{center}

三种尺度里只有最后一种是加法。既然模型只会加法，答案就只剩一个。

\paragraph{替代方案的排除。}\textbf{直接用 $p$ 做线性回归}：$p$ 有界，线性函数必然越界；而且 $p$ 尺度两端极不敏感、中间极敏感，误差结构完全变形。\textbf{直接用 odds}：odds 的自然运算是乘法（式 \eqref{eq:oddsmult}），把它当加法尺度用，等于假设「证据能相加成倍数」，没有任何模型依据。

\begin{dfn}[对数几率]
\begin{equation}
  \mathrm{logit}(p)=\ln\frac{p}{1-p},\qquad p\in(0,1).
  \label{eq:logit}
\end{equation}
它是 $(0,1)$ 与 $\R$ 之间的一一对应，把乘法尺度的几率换成加法尺度的得分，并满足 $\mathrm{logit}(1-p)=-\mathrm{logit}(p)$。
\end{dfn}

几个锚点值得记住，\S6.3 讲校准时要用：

\begin{center}
\small
\begin{tabular}{@{}lccccc@{}}
\toprule
\rowcolor{navy!7}
\tabhead $p$ & $0.01$ & $0.1$ & $0.5$ & $0.9$ & $0.99$\\
\midrule
$\mathrm{logit}(p)$ & $-4.60$ & $-2.20$ & $0$ & $2.20$ & $4.60$\\
\bottomrule
\end{tabular}
\end{center}

这张表说明了两件事。其一，$p$ 稍微离开 $0$ 或 $1$，对数几率就迅速跑远——概率尺度对两端极不敏感，对数几率尺度对中间最敏感。其二，$p$ 与 $1-p$ 对应的两个数关于原点对称，这正是「等距」的含义；换成几率看，就是「互为倒数」，换成对数看，就是「互相变号」。

刻度定好了，下一步是把它翻回概率——这正是 sigmoid 的来历（\S2.2）。至于「凭什么假设这个刻度是线性的」，留给 \S2.3。

\subsection{对数几率的反函数}

现在把假设提出来。既然 $z=\bw\tp\bx$ 可以取任意实数，概率必须落在 $(0,1)$，最自然的一条路是：\textbf{假设对数几率是线性的}，
\begin{equation}
  \ln\frac{p(y=1\mid\bx)}{1-p(y=1\mid\bx)}=\bw\tp\bx .
  \label{eq:logitlinear}
\end{equation}
把式 \eqref{eq:logitlinear} 反解出来，$p$ 的形式就被\emph{唯一}确定了：

\begin{prop}[sigmoid 是反解，不是选择]
式 \eqref{eq:logitlinear} 等价于
\begin{equation}
  p(y=1\mid\bx)=\sigma(\bw\tp\bx),\qquad
  \sigma(z)=\frac{1}{1+e^{-z}} .
  \label{eq:sigmoid}
\end{equation}
并且
\begin{equation}
  \sigma'(z)=\sigma(z)\bigl(1-\sigma(z)\bigr),\qquad
  \sigma(-z)=1-\sigma(z).
  \label{eq:sigmoidid}
\end{equation}
\end{prop}

\begin{proof}
由 $z=\ln\frac{p}{1-p}$ 得 $e^{z}=\frac{p}{1-p}$，即 $p=\frac{e^{z}}{1+e^{z}}=\frac{1}{1+e^{-z}}$。
求导：$\sigma'(z)=\frac{e^{-z}}{(1+e^{-z})^{2}}=\sigma(z)(1-\sigma(z))$。
对称性：$\sigma(-z)=\frac{1}{1+e^{z}}=\frac{e^{-z}}{1+e^{-z}}=1-\sigma(z)$。
\end{proof}

这一步值得停一下，因为绝大多数教材把它写成了「我们\emph{选择}用 sigmoid 把线性函数压成概率」。按刚才的推导，这不是选择：\textbf{只要你相信对数几率是线性的，sigmoid 就是被迫出现的形式}。你能选的只有「对数几率是否线性」这一件事，而这件事恰恰是有内容的假设（\S2.3 会说清它的来源）。

式 \eqref{eq:sigmoidid} 的两条恒等式后面反复要用：第一条把 Hessian 变成 $\bW$ 的加权和（\S4.2），第二条把「判 1 的概率」与「判 0 的概率」接上（\S6.1）。

\begin{add}[从「几率是指数」读出来]
把式 \eqref{eq:logitlinear} 两边取指数，得 $\dfrac{p}{1-p}=e^{\bw\tp\bx}$：\textbf{几率随特征按指数变化}。
每把一个特征 $x^{(j)}$ 增加 $1$ 个单位，几率就乘以 $e^{w^{(j)}}$——这句话是逻辑回归最好用的解释，比「概率加了多少」干净得多，因为概率的加法依赖起点，几率只会乘。
\end{add}
\subsection{线性假设的来源}

式 \eqref{eq:logitlinear} 的假设不是凭空来的。上一讲 §6.1 已经算过一次同样的事：当每一类都是高斯、且共享同一个协方差矩阵时，判别函数是
\begin{equation}
  \delta_c(\bx)=\bx\tp\bSigma^{-1}\bmu_c-\tfrac12\bmu_c\tp\bSigma^{-1}\bmu_c+\log\phi_c ,
  \label{eq:deltac}
\end{equation}
它对 $\bx$ 线性。两类的后验之比就是两个 $\delta$ 之差，于是
\begin{equation}
  \log\frac{p(y=1\mid\bx)}{p(y=2\mid\bx)}
  =\delta_1(\bx)-\delta_2(\bx)
  =(\bmu_1-\bmu_2)\tp\bSigma^{-1}\bx+\text{const}.
  \label{eq:gdaLogit}
\end{equation}
\textbf{共享协方差的高斯假设，恰好生产出一条线性的对数几率。}这条式子在上一讲 §6.3 已经写出，本讲直接引用，不重推。

要做的事只有一件：把类标换过来。上一讲用 $\{1,2\}$，本讲用 $\{0,1\}$，对应关系是
\[
  \text{本讲的类 }1=\text{上一讲的类 }1,\qquad
  \text{本讲的类 }0=\text{上一讲的类 }2 .
\]
换完之后，式 \eqref{eq:gdaLogit} 就是本讲的式 \eqref{eq:logitlinear}，只需把系数记成 $\bw$。

\begin{warn}[一个数值上的坑]
第二讲例题的 $\hat\bw=(0.622,\,0.783)\tp$ 指向它的类 2，也就是本讲的类 0。若把这个数直接搬来当本讲的 $\bw$，$z$ 的符号正好反了。本讲一律用 $\bmu$ 表达，用不到这个数；要借用请先反号。
\end{warn}

现在可以看清这条假设的分量了。式 \eqref{eq:gdaLogit} 的推导用到了「高斯 + 共享协方差」这一整套假设，但\emph{结论}只用到其中一条：对数几率是线性的。逻辑回归做的事，就是\textbf{把前面那些假设全丢掉，只留下这一条}。

这正是判别式与生成式的分界点。生成式建模 $p(\bx\mid y)$，先假设每一类长什么样，再由贝叶斯公式算出后验；判别式不问 $p(\bx\mid y)$，直接对后验的形状下注。两者的差别不在数学难度，而在\emph{赌什么}：
\begin{itemize}
  \item 赌「每一类是高斯」，赢面是样本效率——只要赌对，用很少的样本就能估准（\S7.2）；
  \item 赌「对数几率线性」，赢面是稳健——类条件分布长什么样都无所谓，只要后验的形状对，估计就成立。
\end{itemize}
注意后一条比前一条\emph{弱}：高斯假设成立时，logit 必然线性；反过来不成立。所以逻辑回归的适用范围严格包含 LDA 与 GDA 的高斯情形，代价是假设更弱、同样样本量下估得更糙。

\subsection{链接函数的比较}

\begin{add}[还有哪些「压」法]
任何单调函数 $g:(0,1)\to\R$ 都能当\textbf{链接函数}，把概率接到线性得分上：$g(p)=\bw\tp\bx$。常见的两个替代品是
\[
  \text{probit}:\ g(p)=\Phi^{-1}(p),\qquad
  \text{tanh 形式}:\ p=\tfrac12\bigl(1+\tanh(z/2)\bigr),
\]
其中 $\Phi$ 是标准正态的分布函数。三者形状接近，差别在尾部：logit 的尾部是\textbf{线性}的（$z\to-\infty$ 时 $\sigma(z)\approx e^{z}$），probit 的尾部更薄（$\Phi(-z)\approx\varphi(z)/z$，按 $e^{-z^{2}/2}$ 衰减）。这意味着同样的极端样本，probit 给出更小的概率——它对「离群」更不宽容。

具体选哪个，实践中差别不大；教科书偏爱 logit，有一个干净的理由（见下面的命题）。
\end{add}

\begin{prop}[logit 是伯努利分布的自然参数]
把伯努利分布写成指数族形式
\begin{equation}
  p(y)=p^{y}(1-p)^{1-y}=\exp\Bigl\{y\ln\frac{p}{1-p}+\ln(1-p)\Bigr\},
\end{equation}
则自然参数恰为 $\eta=\ln\frac{p}{1-p}=\mathrm{logit}(p)$，对数配分函数为 $A(\eta)=\ln(1+e^{\eta})$，
且 $A'(\eta)=\sigma(\eta)=p$。
\end{prop}

\begin{proof}
按 $p(y)=\exp\{y\eta-A(\eta)\}$ 对照：$\eta=\mathrm{logit}(p)$，$A(\eta)=-\ln(1-p)=\ln(1+e^{\eta})$。
求导 $A'(\eta)=\dfrac{e^{\eta}}{1+e^{\eta}}=\sigma(\eta)$；而指数族里 $A'(\eta)$ 正是充分统计量 $y$ 的期望，即 $p$。
\end{proof}

这条命题解释了为什么 GLM 的教科书默认用 logit：\textbf{线性预测子取自然参数，是最省事、也最有理论保证的一档}。用统计的话说，$\bw\tp\bx$ 直接等于自然参数时，似然是「标准形式」，得分函数与信息矩阵都有干净的表达式；换成 probit，就没有这层结构了。

\begin{keybox}[本节结论]
\textbf{假设只有一个：对数几率是线性的。}它可能来自高斯假设（上一讲已证），也可以独立地作为出发点。一旦接受它，sigmoid 就是反解出来的唯一形式，$p=0.7$ 也就有了刻度——这里的 $0.7$ 是 $\mathrm{logit}(0.7)=0.847$ 那个位置，与其他模型的 $0.7$ 可以比。
\end{keybox}

% ============================================================
%  第 3 节
% ============================================================

\section{伯努利似然与交叉熵}

\subsection{伯努利似然}

形式定了，接下来是「哪组参数最好」。先把标签的约定钉死。

\begin{warn}[三套标签不能混]
第二讲的 LDA 用类标 $\{1,\dots,K\}$ 并靠类内计数求和；感知机用 $\{+1,-1\}$，因为它要把「错分」写成 $-y\,\bw\tp\bx>0$。本讲用 $\{0,1\}$：只有 $\{0,1\}$ 才能把伯努利分布写成 $p^{y}(1-p)^{1-y}$ 这种紧凑形式。需要 $\pm1$ 时，一律显式写 $s_i=2y_i-1$，不要中途改约定。
\end{warn}

\begin{center}
\small
\begin{tabular}{@{}llll@{}}
\toprule
\rowcolor{navy!7}
\tabhead 讲次 & \tabhead 标签 & \tabhead 典型写法 & \tabhead 用途\\
\midrule
第二讲 & $1,\dots,K$ & $\sum_{i:y_i=c}$ & 类均值、散布矩阵\\
第二讲 & $\pm1$ & $-y_i\,\bw\tp\bx_i>0$ & 感知机的错分判据\\
本讲 & $0,1$ & $p^{y_i}(1-p)^{1-y_i}$ & 伯努利似然\\
\bottomrule
\end{tabular}
\end{center}

在 $\{0,1\}$ 下，单个样本的似然就是一次伯努利试验：
\begin{equation}
  p(y_i\mid\bx_i;\bw)=p_i^{y_i}(1-p_i)^{1-y_i},\qquad p_i=\sigma(\bw\tp\bx_i).
  \label{eq:like}
\end{equation}
$y_i=1$ 时它退化成 $p_i$，$y_i=0$ 时退化成 $1-p_i$——两个分支被一个式子装下了。这正是 $\{0,1\}$ 记号的好处。

假设 $N$ 个样本独立，整份数据的似然是连乘：
\begin{equation}
  \mathcal{L}(\bw)=\prod_{i=1}^{N}p_i^{y_i}(1-p_i)^{1-y_i}.
  \label{eq:likeN}
\end{equation}
极大似然就是找出使 $\mathcal{L}$ 最大的 $\bw$。有两条理由必须先取对数：
\begin{enumerate}[label=(\arabic*),leftmargin=2.0em]
  \item \textbf{防下溢}：$N=10^{4}$ 个 $(0,1)$ 之间的数连乘，double 也存不下，早就是 $0$ 了；取对数变成求和，量级线性增长。
  \item \textbf{保序}：$\ln$ 单调递增，$\argmax\mathcal{L}=\argmax\ln\mathcal{L}$，最大值点不变。
  \item \textbf{与 $\exp$ 相消}：$p_i=\sigma(z_i)$ 里有一个 $\exp$，取对数后正好把它解开，得到 \S3.3 那个极简的梯度。
\end{enumerate}

\subsection{交叉熵与 KL 散度}

\begin{thm}[似然、交叉熵、KL 是同一件事]
记 $\ell_i=\ln p(y_i\mid\bx_i;\bw)$，$L=-\frac1N\sum_i\ell_i$，则
\begin{equation}
  \argmax_{\bw}\ \mathcal{L}(\bw)
  =\argmin_{\bw}\ L(\bw)
  =\argmin_{\bw}\ \mathrm{KL}\bigl(\hat p_{\text{data}}\,\big\Vert\,p_{\text{model}}\bigr),
  \label{eq:equiv}
\end{equation}
其中 $L$ 就是平均交叉熵，$\hat p_{\text{data}}$ 是经验分布。
\end{thm}

\begin{proof}
第一项等号来自 $\ln$ 单调。展开 $L$：
\[
  L=-\frac1N\sum_{i}\Bigl[y_i\ln p_i+(1-y_i)\ln(1-p_i)\Bigr]
  =-\frac1N\sum_i\sum_{c\in\{0,1\}}\mathbb{I}(y_i=c)\ln p(y=c\mid\bx_i;\bw).
\]
把经验分布记成 $\hat p_{\text{data}}(y\mid\bx_i)=\mathbb{I}(y=y_i)$，上式即 $-\mathbb{E}_{\hat p_{\text{data}}}[\ln p_{\text{model}}]$，也就是交叉熵。而
\[
  \mathrm{KL}\bigl(\hat p_{\text{data}}\Vert p_{\text{model}}\bigr)
  =\underbrace{\mathbb{E}_{\hat p_{\text{data}}}\bigl[\ln\hat p_{\text{data}}\bigr]}_{\text{与 }\bw\text{ 无关}}
  -\underbrace{\mathbb{E}_{\hat p_{\text{data}}}\bigl[\ln p_{\text{model}}\bigr]}_{=-L},
\]
第一项是经验分布的熵，不含参数，所以最小化 KL 与最小化 $L$ 等价。
\end{proof}

这个定理值得记住的形状是：\textbf{「训练集上最像」这句话，有三种等价说法}。说「似然最大」是在讲概率，「交叉熵最小」是在讲信息，「KL 最小」是在讲两个分布的差异。三者同源，所以后文可以随手换用最方便的那一个。

同时也要记住它\emph{不}保证什么：让 $L$ 最小，只说明模型在\emph{训练集的经验分布}上最接近，既不保证泛化，也不保证概率可以当频率用。后者是 \S6.3 的事。

\begin{add}[为什么不是均方误差]
用平方误差 $L_{\text{MSE}}=\frac1N\sum(y_i-p_i)^2$ 也能训练，但它的梯度里会多出一个因子 $p_i(1-p_i)$：$\partial L_{\text{MSE}}/\partial z_i=2(p_i-y_i)p_i(1-p_i)$。当模型已经错得离谱（$p_i\to0$ 而 $y_i=1$）时，这个因子趋近 $0$，梯度反而消失——\textbf{错得越狠，改得越慢}。交叉熵没有这个问题：它的梯度就是 $p_i-y_i$，错得越狠，改得越快。

这也是「交叉熵 = 对数似然」比「平方误差」更适合分类的根本原因，与凸性一起，构成本讲 \S4 的两个好性质。
\end{add}

\subsection{梯度的形式}

有了 $L$，求梯度只是一次链式法则。

\begin{prop}[梯度是「残差 × 特征」]
对单样本对数似然 $\ell_i=y_i\ln p_i+(1-y_i)\ln(1-p_i)$，
\begin{equation}
  \nabla_{\bw}\ell_i=(y_i-p_i)\,\bx_i,\qquad
  \bg=\nabla L=\frac1N\sum_{i=1}^{N}(p_i-y_i)\,\bx_i .
  \label{eq:grad}
\end{equation}
\end{prop}

\begin{proof}
对 $z_i$ 求导，注意 $p_i=\sigma(z_i)$ 且 $\sigma'=p(1-p)$（式 \eqref{eq:sigmoidid}）：
\[
  \frac{\partial L_i}{\partial z_i}
  =-\Bigl[\frac{y_i}{p_i}p_i(1-p_i)-\frac{1-y_i}{1-p_i}p_i(1-p_i)\Bigr]
  =-y_i(1-p_i)+(1-y_i)p_i
  =p_i-y_i ,
\]
其中 $L_i=-\ell_i$ 是单样本损失；再由 $z_i=\bw\tp\bx_i$ 得 $\partial z_i/\partial\bw=\bx_i$，乘起来即为 \eqref{eq:grad}。
\end{proof}

两个负号要盯住：$\ell$ 是正项似然，所以 $\nabla\ell$ 是 $y_i-p_i$；取负号变成损失，才得到 $p_i-y_i$。写代码时若符号搞反，梯度上升会变成下降，损失会一路增大——这是最常见的一处笔误。

式 \eqref{eq:grad} 的形状值得单独看一眼：它是\textbf{残差乘以特征}，与最小二乘的梯度 $2\sum(\hat y_i-y_i)\bx_i$ 完全同构，只差一个因子 $2$（以及这里的 $\hat y_i$ 换成了 $p_i$）。这个相似不是巧合：两者都是「指数族 + 线性预测子」的产物，也都通向同一种算法——加权最小二乘（\S4.3）。

\begin{warn}[$\pm1$ 标签代进似然会失效]
若把 $y_i\in\{+1,-1\}$ 直接代入 $p^{y}(1-p)^{1-y}$，$y=-1$ 时指数变成 $2$，式子毫无意义。要写 $\pm1$ 版本，必须改成 $\ln p(y_i\mid\bx_i)=-\ln\bigl(1+e^{-y_i z_i}\bigr)$，这是另一种等价写法，\S5.1 判断分离时会用到。
\end{warn}

\begin{keybox}[本节结论]
\textbf{估计准则不是「想出来的」，是似然推出来的。}$L$ 就是平均交叉熵，它的梯度是残差乘特征，形状与最小二乘一致。这条式子同时解释了为什么用交叉熵而不用平方误差：前者错得越狠梯度越大，后者恰好相反。
\end{keybox}

% ============================================================
%  第 4 节
% ============================================================

\section{凸性与重加权最小二乘}

\subsection{无闭式解的原因}

令梯度为零：
\begin{equation}
  \bg=\frac1N\sum_{i=1}^{N}\bigl(\sigma(\bw\tp\bx_i)-y_i\bigr)\bx_i=\bzero .
  \label{eq:firstorder}
\end{equation}
这是 $(D+1)$ 个方程、$(D+1)$ 个未知数，看起来与最小二乘的正规方程没差别。差别在于解不出来。

\textbf{先把理由说准。}常见的错误解释是「因为概率是非线性的，所以没有闭式解」——这句话站不住：最小二乘的正规方程 $\bX\tp\bX\bw=\bX\tp\by$ 对 $\bw$ 是\textbf{线性}的，所以能直接解；而它的残差平方本身关于 $\bw$ 是二次的，求导后自然降到一次。非线性并不妨碍求解，多项式方程也有公式解（至少低次有）。

真正的原因是：$\sigma$ 把 $\bw$ 包在指数里，式 \eqref{eq:firstorder} 成为一个\textbf{超越方程组}——$\bw$ 同时出现在指数上和系数上，没有任何代数变形能把它单独解出来。这跟 $x e^{x}=1$ 解不出初等表达式是同一类困难。

所以逻辑回归的 MLE 只能\emph{迭代}求。而迭代之前，先要确认这个方程的解存在、唯一、且值得追——这是 \S4.2 的事。

\subsection{Hessian 与凸性}

凸性看二阶导，所以先把它算出来。起点是 \S3.3 的单样本导数 $\partial L_i/\partial z_i=p_i-y_i$，再对 $z_i$ 求一次：
\begin{equation}
  \frac{\partial^{2}L_i}{\partial z_i^{2}}
  =\frac{\partial\,(p_i-y_i)}{\partial z_i}
  =\frac{\partial p_i}{\partial z_i}
  =p_i(1-p_i).
  \label{eq:d2l}
\end{equation}
这一步只用 $\sigma'=\sigma(1-\sigma)$。但 $p_i(1-p_i)$ 不是凑出来的系数，它有两个身份：是伯努利变量的\textbf{方差} $\Var(y_i)$，也是 $\sigma$ 在 $z_i$ 处的\textbf{斜率}。两个身份说的是同一件事——预测越靠近 $0.5$，曲率越大；越靠近 $0$ 或 $1$，曲率越小，$L$ 在那附近几乎是平的。

把 $z_i=\bw\tp\bx_i$ 代回去，链式法则给出
\begin{equation}
  \nabla^{2}L
  =\frac1N\sum_{i=1}^{N}\frac{\partial^{2}L_i}{\partial z_i^{2}}\,\bx_i\bx_i\tp
  =\frac1N\sum_{i=1}^{N}p_i(1-p_i)\,\bx_i\bx_i\tp ,
  \label{eq:hessiansum}
\end{equation}
一个\textbf{加权外积和}：每个样本贡献一个秩一矩阵，权是它的方差。要写成矩阵形式，先给权重起个名字：
\begin{equation}
  \bW=\diag\bigl(p_i(1-p_i)\bigr)\in\R^{N\times N},
  \label{eq:W}
\end{equation}
于是
\begin{equation}
  \bH=\nabla^{2}L=\frac1N\bX\tp\bW\bX .
  \label{eq:hessian}
\end{equation}
这一步没有新的推导，只是把求和改写。$\bX\tp\bW\bX$ 的第 $(j,k)$ 个元素按矩阵乘法是
\[
  \sum_{i=1}^{N}x_i^{(j)}\,p_i(1-p_i)\,x_i^{(k)},
\]
正是 \eqref{eq:hessiansum} 里的第 $(j,k)$ 项。\textbf{所以 $\bW$ 既不是新参数，也不是超参数，它是「每个样本的方差」摞成的对角阵}；\S4.3 的 IRLS 会再遇见它，是同一个东西。

它与参数 $\bw$ 只共享一个字母，含义毫不相干：$\bw\in\R^{D+1}$ 是\textbf{待估}的系数，训练开始前未知；$\bW\in\R^{N\times N}$ 是由\emph{当前}的 $\bw$ 算出来的\textbf{权重}，每迭代一步重算一次。换句话说，$\bw$ 是未知量，$\bW$ 是它的函数。（也不要与第二讲的 $\bW$ 混：那里是投影矩阵 $D\times d$，本讲是权重阵 $N\times N$，同名不同物。）

\begin{thm}[损失严格凸]
在 \eqref{eq:hessian} 的记号下 $\bH\succeq0$；当 $\bX$ 列满秩且 $0<p_i<1$ 对所有 $i$ 成立时，$\bH$ 正定，$L$ 严格凸。此时\textbf{局部极小即全局极小}，极小点唯一。
\end{thm}

\begin{proof}
对任意 $\bv$，
$\bv\tp\bH\bv=\frac1N\sum_i p_i(1-p_i)(\bv\tp\bx_i)^2\ge0$，故半正定。若 $\bX$ 列满秩且所有权重为正，$\bv\tp\bH\bv=0$ 迫使 $\bv\tp\bx_i=0$ 对所有 $i$ 成立，即 $\bX\bv=\bzero$，只有 $\bv=\bzero$，故正定。
\end{proof}

顺带一句，这个 $\bH$ 就是（负的）Fisher 信息矩阵除以 $N$：对 canonical link 的 GLM，观测信息与期望信息相等，而 $\bW$ 的对角元既是方差又是信息量——在指数族里这两件事本来就是同一个数。所以「半正定」不是巧合，只是「方差非负」的矩阵版本。

凸性这件事在分类里并不常见，值得对比一下：感知机的损失 $\sum\max(0,-y_i z_i)$ 分段线性、非凸，解不唯一，初始化和样本顺序都会改变结果；逻辑回归的损失严格凸（在非退化数据上），只要算法收敛，落到的就是同一个点。\textbf{「换一份实现、结果不一样」这种事，在逻辑回归里不该发生}。如果发生了，先怀疑数据（分离，见 \S5），再怀疑代码。

两个前提要记住。$\bX$ 列满秩：$D+1>N$ 时必然不满秩，此时有无穷多组解给出同样的 $L$，需要正则化（\S5.3）。$0<p_i<1$：一旦某个 $p_i$ 精确等于 $0$ 或 $1$（数值饱和或完全分离），对应的权重变成 $0$，$\bH$ 退化——这正是 \S4.4 与 \S5 要处理的两件事。

\subsection{牛顿法与 IRLS}

既然有梯度与 Hessian，最自然的算法是牛顿法：在当前点用二次函数近似目标，跳到二次函数的极小点。
\begin{equation}
  \bw^{(t+1)}=\bw^{(t)}-\bH^{-1}\bg
  =\bw^{(t)}-\bDelta^{(t)},
  \qquad \bDelta^{(t)}=\bH^{-1}\bg .
  \label{eq:newton}
\end{equation}
把 $\bH$ 与 $\bg$ 的具体表达式代进去，会看到一个熟悉的结构：
\begin{equation}
  \bDelta=\Bigl(\tfrac1N\bX\tp\bW\bX\Bigr)^{-1}\tfrac1N\bX\tp(p-y)
  \quad\Longrightarrow\quad
  \bX\tp\bW\bX\,\bDelta=\bX\tp(p-y).
  \label{eq:newtonstep}
\end{equation}
$1/N$ 在两边同时出现，约掉了。

\textbf{这里有一个极容易被读错的地方。}式 \eqref{eq:newtonstep} 是牛顿步，\emph{不是}「以 $\bW$ 为权、拟合 $p-y$」的正规方程。照后一种说法写，右边还得再乘一个 $\bW$：
\begin{equation}
  \bX\tp\bW\bX\,\bDelta=\bX\tp\bW(p-y),
  \label{eq:wrongwls}
\end{equation}
它与 \eqref{eq:newtonstep} 并不等价，而且差别不是「差一个常数倍」。取 \S4.5 例题第 3 步的值代入试算：两种写法给出的增量，逐分量比值是 $0.09$、$-1.68$、$-0.64$，夹角余弦 $-0.94$——连方向都不同。

原因在于 $\bW$ 出现的\emph{次数}不同。在加权最小二乘里，目标函数本身带权，所以 $\bW$ 在正规方程左边出现一次（进了曲率）、右边出现一次（乘在响应上）；在牛顿步里，$\bW$ 只是\emph{曲率}，梯度 $\bg$ 里没有它。同一个 $\bW$ 不能既当曲率、又在右边当残差的权重。

这个错误在第一步还看不出来：$t=0$ 时 $p_i\equiv1/2$，$\bW=\tfrac14\bI$ 是数量阵，\eqref{eq:wrongwls} 与 \eqref{eq:newtonstep} 只差常数因子 $1/4$，看起来「只是步长小了点」。等各点的 $p_i$ 拉开差距、权重不再一样，方向才会歪。

IRLS 的正确读法要把未知量与响应一起换掉：未知量取「下一个点」而不是增量，响应取下面的\textbf{工作响应}：
\begin{equation}
  \bw^{(t+1)}=\bigl(\bX\tp\bW\bX\bigr)^{-1}\bX\tp\bW\,\bz_{\mathrm w},
  \qquad
  \bz_{\mathrm w}=\bX\bw^{(t)}+\bW^{-1}(y-p).
  \label{eq:irls}
\end{equation}
$\bz_{\mathrm w}$ 是 $N$ 维向量，与单样本得分 $z_i$ 不是一回事，下标 $\mathrm w$ 就是为了把两者分开。

工作响应还有一层几何含义：在当前点用切线近似 $p$，要让线性化的预测恰好等于 $y_i$，得分该取多少。由
\[
  y_i\approx p_i+p_i(1-p_i)\bigl(z_i^{\mathrm{new}}-z_i\bigr)
\]
反解即得 $z_i^{\mathrm{new}}=z_i+\dfrac{y_i-p_i}{p_i(1-p_i)}$，写成向量就是 \eqref{eq:irls} 里的 $\bz_{\mathrm w}$。

\begin{prop}[IRLS：每一步都是一次加权最小二乘]
式 \eqref{eq:irls} 的迭代，等价于：以当前权重 $\bW$ 做一次加权最小二乘，\textbf{拟合工作响应 $\bz_{\mathrm w}$}（不是拟合残差 $p-y$），解出的是新的参数值而非增量。权重与响应都随迭代更新，故称\textbf{迭代重加权最小二乘}（IRLS）。
\end{prop}

\begin{proof}
在式 \eqref{eq:newtonstep} 中解出 $\bw^{(t+1)}=\bw^{(t)}-\bDelta$，代入 $\bDelta$：
\[
  \bw^{(t+1)}=\bw^{(t)}-\bigl(\bX\tp\bW\bX\bigr)^{-1}\bX\tp(p-y)
  =\bigl(\bX\tp\bW\bX\bigr)^{-1}\bX\tp\bW\Bigl[\bX\bw^{(t)}+\bW^{-1}(y-p)\Bigr],
\]
最后一步用了 $\bX\tp\bW\cdot\bW^{-1}(y-p)=\bX\tp(p-y)$，并把 $\bw^{(t)}$ 塞进括号。
\end{proof}

换句话说，\textbf{$\bW$ 并没有消失，它与藏在工作响应里的 $\bW^{-1}$ 相消了}——这正是 \eqref{eq:newtonstep} 右边看不见 $\bW$ 的原因。记住这条，就不会再把式 \eqref{eq:newtonstep} 读成「以 $\bW$ 为权、响应是 $p-y$ 的加权最小二乘」。

这条命题把第一讲的机器原样搬了过来：\textbf{加权最小二乘你早就会算}，逻辑回归的训练不过是「算一次、更新权重、再算一次」。这也解释了为什么早期的统计软件能直接用回归模块拟合逻辑回归——只要允许传权重。

与梯度下降对比一下，两者的差别在收敛速度与每步代价：

\begin{center}
\small
\begin{tabular}{@{}lccc@{}}
\toprule
\rowcolor{navy!7}
\tabhead 方法 & \tabhead 每步代价 & \tabhead 收敛阶 & \tabhead 是否需调步长\\
\midrule
梯度下降 & $O(ND)$ & 线性（误差每步乘常数） & 需要\\
牛顿法 / IRLS & $O(ND^{2}+D^{3})$ & 二次（误差平方） & 不需要\\
\bottomrule
\end{tabular}
\end{center}

「二次收敛」的意思是：一旦靠近解，有效数字每步翻倍。\S4.5 的例题会看到，从零出发七步就到底了。代价是每步要解一个 $(D+1)$ 阶线性方程组——$D$ 上千时这一步就很贵，于是有了拟牛顿（L-BFGS）与小批量随机梯度下降。\textbf{中小规模用牛顿，大数据用一阶方法}，这是实践中的分界线。

\subsection{数值实现要点}

公式写完，能跑起来是另一回事。四条经验，每条都对应一个具体的坑。

\paragraph{softplus 稳定式。}交叉熵的两个分量都是 softplus：
\begin{equation}
  -\ln\sigma(z)=\ln\bigl(1+e^{-z}\bigr)=\softplus(-z),\qquad
  -\ln\bigl(1-\sigma(z)\bigr)=\ln\bigl(1+e^{z}\bigr)=\softplus(z),
\end{equation}
而 softplus 的稳定式是
\begin{equation}
  \softplus(u)=\max(u,0)+\ln\bigl(1+e^{-\abs{u}}\bigr).
  \label{eq:softplus}
\end{equation}
\textbf{恒等式右边是 $\softplus(u)$，不是 $\log(1+e^{-u})$ 配 $\max(u,0)$}——这两者差一个符号。若照着 $\max(z,0)$ 去算 $\ln(1+e^{-z})$，$z=2$ 时会得到 $2.127$，而正确答案是 $0.127$。写法上把自变量写成 $u$ 就是为了避免这种错配。

\paragraph{工作响应的实现。}$\bW$ 的对角元是 $p_i(1-p_i)$，当 $p_i\to0$ 或 $1$ 时它们趋近 $0$，$\bW^{-1}$ 的元素会爆炸。$p\to0/1$ 的情形在训练后期很常见，此时 $\bz_{\mathrm w}$ 的表达式只是\emph{推导}上的方便，\emph{实现}上应直接解加权正规方程
\begin{equation}
  \bigl(\bX\tp\bW\bX\bigr)\,\bDelta=\bX\tp(p-y),
\end{equation}
或用 QR、Cholesky 分解。$\bDelta$ 是牛顿增量，不是 $\delta_c$。

\paragraph{停止准则。}用「目标函数差小于容差」当判据在逻辑回归上很危险，本讲例题就有两个现成的反例：其一，$p$ 一旦精确饱和到 $1$，朴素式 $-\frac1N\sum[y\ln p+(1-y)\ln(1-p)]$ 会直接返回 \texttt{nan}（$0\cdot\ln0$），只有式 \eqref{eq:softplus} 才拿得到有限值；其二，在完全分离的数据上，$L$ 与 $\norm{\bg}$ 会一路掉到机器精度，而 $\norm{\bw}$ 还在线性增长（\S5.2）。用 $L$ 作判据，会在参数完全没收敛时宣布成功。

\paragraph{特征标准化与条件数。}$\bH=\frac1N\bX\tp\bW\bX$ 的条件数随特征尺度的差异平方增长。两个特征一个以「米」计、一个以「毫米」计，$\bH$ 的条件数就是 $10^{6}$ 量级，牛顿步会在两个方向上跳得极不均匀。标准化之后再看 $\bH$ 是否近奇异；若近奇异，加一个小量 $\epsilon\bI$，或改用 L-BFGS。

\subsection{数值例题}

\begin{exm}[二维六点，手算可复现]
数据如下（$y=1$ 为前三点）。表中只列前两维；按本讲约定，$\bx_i$ 的末位恒为 $1$（截距位），不是第三个特征：
\begin{center}
\small
\begin{tabular}{@{}lcccccc@{}}
\toprule
\rowcolor{navy!7}
\tabhead $i$ & 1 & 2 & 3 & 4 & 5 & 6\\
\midrule
$\bx_i$（前两维） & $(1,\,1)$ & $(1.5,\,1.2)$ & $(2.7,\,2.4)$ & $(2.6,\,2.1)$ & $(3.2,\,2.6)$ & $(2.1,\,3.1)$\\
$y_i$ & $1$ & $1$ & $1$ & $0$ & $0$ & $0$\\
\bottomrule
\end{tabular}
\end{center}
这六个点\textbf{不可分}：第 3 点落在类 0 三点的凸包内（重心坐标约 $(0.553,\,0.294,\,0.153)$，严格在内部），所以任何直线都无法把它们完全分开，MLE 存在（\S5.1 会解释这条推理）。

从 $\bw=\bzero$ 出发做牛顿迭代（等价 IRLS），七步收敛：
\begin{center}
\small
\begin{tabular}{@{}rcccccc@{}}
\toprule
\rowcolor{navy!7}
\tabhead $t$ & \tabhead $w_1$ & \tabhead $w_2$ & \tabhead $b$ & \tabhead $\norm{\bg^{(t)}}$ & \tabhead $L^{(t)}$\\
\midrule
0 & $0$ & $0$ & $0$ & $3.489\times10^{-1}$ & $0.6931$\\
1 & $-0.4263$ & $-1.5911$ & $4.2192$ & $2.426\times10^{-2}$ & $0.4044$\\
2 & $-0.6490$ & $-2.2358$ & $6.2561$ & $1.457\times10^{-2}$ & $0.3754$\\
3 & $-0.8977$ & $-2.5206$ & $7.6026$ & $4.494\times10^{-3}$ & $0.3701$\\
4 & $-1.0154$ & $-2.6167$ & $8.1512$ & $4.084\times10^{-4}$ & $0.3696$\\
5 & $-1.0305$ & $-2.6287$ & $8.2205$ & $4.893\times10^{-6}$ & $0.3696$\\
6 & $-1.0307$ & $-2.6288$ & $8.2214$ & $8.596\times10^{-10}$ & $0.3696$\\
7 & $-1.0307$ & $-2.6288$ & $8.2214$ & $6.029\times10^{-16}$ & $0.3696$\\
\bottomrule
\end{tabular}
\end{center}
解为 $\hat{\bw}=(-1.0307,\,-2.6288,\,8.2214)$，平均交叉熵 $L=0.3696$。注意 $\norm{\bg}$ 从 $0.35$ 降到 $10^{-16}$，而 $L$ 从第 4 步起就稳定在四位小数不动了——\textbf{用 $L$ 的变化量作判据，第 5 步就会误判收敛}。

各点概率为
\[
  \bp=(0.9897,\ 0.9713,\ 0.2951,\ 0.5053,\ 0.1288,\ 0.1098).
\]
第 3 点（$p=0.2951$，真值 1）与第 4 点（$p=0.5053$，真值 0）值得单独看：前者是「落错侧」的点，模型给它三成把握，而不是硬把它圈回去；后者\emph{刚刚}越过 $0.5$，因此被判为 1，也是错分。

于是交叉熵最优解共错两个点。但下面这条边界只错一个：
\[
  x+y=3\quad(\text{判 1 当 } x+y<3),
\]
它把第 1、2 点判 1（对），第 3 点判 0（错），第 4、5、6 点判 0（对）。\textbf{错分更少的边界，交叉熵反而更大}——这不是逻辑回归的缺陷，而是它的目标就写在式 \eqref{eq:equiv} 里：优化的是概率，不是错分数。想要错分最少，得换损失（合页、0-1）或改阈值（\S6.4）。
\end{exm}

这个例题同时是封面插图的坐标来源：六个点、三条等概率线（$p=0.1,\,0.5,\,0.9$）都从这一组参数算出，正文数字与图来自同一次计算。
% ============================================================
%  第 5 节
% ============================================================

\section{完全分离与正则化}

\subsection{分离的两种形态}

上一节末尾说「不可分，所以 MLE 存在」。这句话反过来问就是：\emph{可分的时候会怎样？}答案不体面——解不存在。

先把「可分」说准。它有两种强度，实践里都遇到，后果略有差别。

\begin{dfn}[完全分离与拟完全分离]
用 $s_i=2y_i-1\in\{\pm1\}$ 记符号化标签，$z_i=\bw\tp\bx_i$。
\begin{enumerate}[label=(\roman*),leftmargin=2.2em]
  \item \textbf{完全分离}：存在 $\bw\ne\bzero$，使 $s_i z_i>0$ 对所有 $i$ \emph{严格}成立；
  \item \textbf{拟完全分离}：存在 $\bw\ne\bzero$，使 $s_i z_i\ge0$ 对所有 $i$ 成立，且\textbf{至少有一个取等号}。
\end{enumerate}
两条都必须要求 $\bw\ne\bzero$：取零向量时 $s_iz_i\equiv0$，(ii) 对\emph{任何}数据集都成立，判据就废了。
\end{dfn}

几何上，这两条对应两类点集的凸包处于什么位置：

\begin{center}
\small
\begin{tabular}{@{}lll@{}}
\toprule
\rowcolor{navy!7}
\tabhead 凸包关系 & \tabhead 分离类型 & \tabhead MLE\\
\midrule
内部相交 & 不可分 & 存在（本讲例题属于这一档）\\
仅在边界相交／相切 & 拟完全分离 & 不存在\\
不相交 & 完全分离 & 不存在\\
\bottomrule
\end{tabular}
\end{center}

「不相交 $\iff$ 完全分离」是凸集分离定理的直接推论：两个紧凸集不相交时，必有一张超平面把它们严格分开。这条判据对任意维成立，二维只是便于把凸包画出来。

拟完全分离值得多说一句，因为它在真实数据里比完全分离更常见。典型来源是\textbf{某个离散特征的一个取值 100\% 属于同一类}：比如「是否 VIP」在正类里全为 1，在负类里全为 0。这时只要让 VIP 对应的系数趋于 $+\infty$，正类的 $z$ 就全部推向 $+\infty$；而负类里那些 VIP $=0$ 的样本，$z$ 由其他特征决定，可能正好落在边界上——于是「至少一个取等号」。这类数据在业务表里到处都是，而且往往「看起来很正常」，直到你发现系数怎么都收敛不了。

\begin{add}[判据既不是万能的，也不是必需的]
Albert \& Anderson (1984) 给出了 MLE 存在性的完整分类：\emph{不可分}时 MLE 存在且唯一；\emph{完全分离或拟完全分离}时不存在。这是充要条件，不只是充分条件。

实践里通常不真的去算凸包，而是看三个信号：迭代次数打满仍不收敛、$\norm{\bw}$ 一路增长、某些样本的 $p_i$ 精确等于 $0$ 或 $1$。三者同时出现，几乎可以肯定就是分离。
\end{add}

\subsection{权重的发散}

\begin{exm}[删掉一个内点，解就跑了]
取 \S4.5 的六个点，删掉第 3 点（那个落在类 0 凸包内的点）。剩下五点完全可分：类 1 是 $(1,1)$、$(1.5,1.2)$，类 0 是 $(2.6,2.1)$、$(3.2,2.6)$、$(2.1,3.1)$，一条 $x+y=3.6$ 就能切开。

仍从 $\bw=\bzero$ 出发做牛顿迭代，观感与上一例完全不同：每一步 $\norm{\bw}$ 增加约 $3.3\to4.0$，几乎没有衰减。第 27 步时
\[
  \bw=(-33.7,\,-18.5,\,99.6),\qquad \norm{\bw}=\sqrt{33.7^{2}+18.5^{2}+99.6^{2}}=106.7,
\]
而 $\norm{\bg}$ 已降到 $10^{-12}$ 以下。参数还在跑，梯度已经看不出问题——这正是最危险的地方。

注意 $106.7$ 这个长度里，截距 $99.6$ 占了大头：把末位摘掉，两维权重只有 $\sqrt{33.7^{2}+18.5^{2}}=38.4$。换句话说，发散最凶的方向其实是截距，不是特征权重。

顺带一个实现上的细节：到第 22 步左右，$p$ 已精确饱和到 $1$。此时朴素式 $-\frac1N\sum[y\ln p+(1-y)\ln(1-p)]$ 会算出 \texttt{nan}（$0\cdot\ln0$），只有 \S4.4 的 softplus 稳定式才拿得到有限值。所以「$\norm{\bg}$ 小」和「$L$ 小」在分离数据上都是假信号。
\end{exm}

\begin{figure}[htbp]
\centering
\begin{tikzpicture}
\begin{axis}[
  name=main,
  width=8.0cm,height=5.6cm,
  xlabel={迭代步 $t$},ylabel={$\norm{\bw^{(t)}}$},
  xmin=0,xmax=28,ymin=0,ymax=115,
  axis lines=left,axis line style={steel!60,line width=0.5pt},
  tick style={steel!60},tick label style={font=\scriptsize,text=steel},
  label style={font=\scriptsize,text=steel},
  legend style={font=\scriptsize,draw=steel!30,fill=white,at={(0.03,0.97)},anchor=north west},
  grid=both,grid style={steel!12,line width=0.3pt},
]
\addplot[navy,line width=0.9pt,mark=none] coordinates {
  (0,0)(1,5.1)(2,8.4)(3,11.9)(4,15.6)(5,19.4)(6,23.3)(7,27.2)(8,31.1)(9,35.1)
  (10,39.1)(11,43.1)(12,47.0)(13,51.0)(14,55.0)(15,59.0)(16,63.0)(17,66.9)
  (18,70.9)(19,74.9)(20,78.9)(21,82.9)(22,86.8)(23,90.8)(24,94.8)(25,98.8)
  (26,102.8)(27,106.7)};
\addlegendentry{完全可分}
\addplot[gold,line width=1.0pt,mark=none] coordinates {
  (0,0)(1,4.5)(2,6.7)(3,8.1)(4,8.6)(5,8.7)(6,8.7)(7,8.7)(8,8.7)};
\addlegendentry{不可分}
\end{axis}
\begin{axis}[
  at={($(main.east)+(0.55cm,0.55cm)$)}, anchor=west,
  width=4.3cm,height=3.3cm,
  xmin=0,xmax=8.4,ymin=0,ymax=10.5,
  title={\scriptsize 前 8 步放大},
  title style={text=steel,font=\scriptsize},
  axis lines=left,axis line style={steel!60,line width=0.5pt},
  tick style={steel!60},tick label style={font=\tiny,text=steel},
  grid=both,grid style={steel!12,line width=0.3pt},
]
\addplot[gold,line width=1.0pt,mark=*,mark size=1.0pt] coordinates {
  (0,0)(1,4.5)(2,6.7)(3,8.1)(4,8.6)(5,8.7)(6,8.7)(7,8.7)(8,8.7)};
\addplot[navy,line width=0.9pt,mark=none] coordinates {(0,0)(1,5.1)(2,8.4)(3,11.9)};
\draw[steel!45,line width=0.4pt,densely dashed] (axis cs:0,8.69) -- (axis cs:8.4,8.69);
\node[font=\tiny,text=steel,anchor=south west] at (axis cs:1.6,8.9) {$8.69$};
\end{axis}
\end{tikzpicture}
\caption{两种数据的 $\norm{\bw^{(t)}}$。可分那条是\textbf{直线增长}，不是指数——牛顿法每步加上一个几乎恒定的向量；不可分那条几步就贴到 $8.69$ 附近不再动（右侧是前八步的放大）。判断「有没有解」，看的就是这两条曲线的分岔。}
\label{fig:growth}
\end{figure}

\begin{prop}[可分时 MLE 不存在]
若数据完全分离，则 $\sup_{\bw}\mathcal{L}=1$，但不存在有限的 $\bw$ 使 $\mathcal{L}=1$。
\end{prop}

\begin{proof}
设 $\bw_0$ 给出分离。对 $t>0$，取 $\bw=t\bw_0$，则 $s_iz_i=t\,s_i\bw_0\tp\bx_i\to+\infty$（$t\to\infty$），于是每个 $p_i\to y_i$，$\mathcal{L}\to1$。但任何有限 $\bw$ 都有某个 $s_iz_i$ 有限，相应 $p_i\ne y_i$，故 $\mathcal{L}<1$。
\end{proof}

于是最优解「在无穷远处」：程序找不到它，只能看到参数一路变大。这就是「发散」的确切含义——不是数值溢出，而是目标函数的上确界无法取到。

\begin{warn}[这不是代码 bug]
新手看到「$\norm{\bw}$ 每轮变大、\texttt{max\_iter} 报收敛失败」，第一反应是调大迭代次数或换求解器。方向反了：\textbf{数据里没有有限解，跑多久都不会有}。正确处置是 \S5.3--\S5.5 的三条路。
\end{warn}

\subsection{L2 正则与 MAP}

最直接的补救是给目标加一个惩罚：
\begin{equation}
  \min_{\bw}\ L(\bw)+\frac{\lambda}{2}\sum_{j=1}^{D}\bigl(w^{(j)}\bigr)^{2},\qquad \lambda>0 .
  \label{eq:l2}
\end{equation}

\begin{thm}[L2 正则等价于高斯先验的 MAP]
在式 \eqref{eq:l2} 中，若对前 $D$ 个分量取各分量独立、方差为 $1/\lambda$ 的零均值高斯先验，对末位（截距 $b=w^{(D+1)}$）取平坦先验，则最小化 \eqref{eq:l2} 与最大化后验 $p(\bw\mid\text{数据})$ 是同一件事。
\end{thm}

\begin{proof}
后验的对数为 $\ln p(\text{数据}\mid\bw)+\ln p(\bw)+C$。第一项是 $N\ln\mathcal{L}=-NL$；第二项在独立高斯先验下为 $-\frac{\lambda}{2}\sum_{j=1}^{D}\bigl(w^{(j)}\bigr)^{2}+C'$（截距位的平坦先验只贡献常数）。两边除以 $N$ 并去掉常数，即得 \eqref{eq:l2}。
\end{proof}

这个等价把第一讲的岭回归接了过来：\textbf{同一个惩罚项，换了个似然}。第一讲是「高斯噪声 + 高斯先验」，这里是「伯努利似然 + 高斯先验」。所以岭回归的那些直觉可以照搬：$\lambda$ 越大，系数越往零拉，方差减、偏差增。

但 \eqref{eq:l2} 多带来一个第一讲没有的好处：\textbf{解一定存在，而且唯一}。原因很直接：加上 $\frac{\lambda}{2}\sum_{j\le D}(w^{(j)})^{2}$ 之后，目标函数关于 $\bw$ 是强制的（coercive，$\norm{\bw}\to\infty$ 时目标 $\to\infty$），所以极小点不会跑到无穷远；而 Hessian 变成
\begin{equation}
  \bH_{\lambda}=\frac1N\bX\tp\bW\bX+\lambda\begin{pmatrix}\bI_{D}&\\&\mathbf{0}\end{pmatrix},
\end{equation}
它在前 $D$ 维上恒正定，与 $\bX$ 是否列满秩无关。于是「$D>N$ 时解不唯一」与「可分时无解」两个毛病一起治好。

注意先验只加在前 $D$ 个分量上、截距位用平坦先验，这不是笔误：截距代表「基准几率」，对它加惩罚等于强行把基准拉到 $1/1$，通常没有道理。代价是截距仍可能发散（例如数据只有一类时，把它推向无穷就能让损失趋零），但前 $D$ 维被压住了，预测不会跟着一起炸。

\subsection{L1 与稀疏解}

把 $\norm{\bw}^{2}$ 换成 $\norm{\bw}_{1}=\sum_j\abs{w^{(j)}}$，就得到 L1 正则（Lasso）：
\begin{equation}
  \min_{\bw,b}\ L(\bw,b)+\lambda\norm{\bw}_{1}.
\end{equation}
几何直觉与第一讲完全相同：L1 的等高线是菱形，与似然等高线相切的位置容易落在坐标轴上，于是部分系数被压到\textbf{恰好为零}——这是一条自动的特征选择。

但这里有一个与算法直接相关的后果：$\norm{\bw}_{1}$ 在原点不可导，\textbf{牛顿法失效}（Hessian 根本不存在），坐标下降与近端梯度（proximal gradient）才是对的方法。近端梯度的一步是「先做梯度下降，再做软阈值」：
\begin{equation}
  w^{(j)}\leftarrow \mathcal{S}_{\lambda\eta}\Bigl(w^{(j)}-\eta\,\partial_{w^{(j)}}L\Bigr),
  \qquad
  \mathcal{S}_{\kappa}(u)=\sign(u)\max\bigl(\abs{u}-\kappa,\,0\bigr).
\end{equation}
软阈值把小的分量直接归零，这就是稀疏性的来源。

一句提醒：\textbf{L1 的稀疏性依赖特征尺度}。若某个特征以「毫米」计、另一个以「米」计，惩罚对前者更狠，选出来的特征会失真。所以用 L1 之前必须标准化——这条在 \S7.3 的清单里还会出现。

\subsection{早停与迭代上限}

第三条路是把「迭代多少步」当作正则：训练到第 $T$ 步就停，$\norm{\bw}$ 自然被限制住，解不会跑到无穷远。

这个说法在大方向上对，但\textbf{有个重要限定}：\emph{早停 $\approx$ L2} 这个结论是在\textbf{梯度下降、步长足够小}的条件下推出来的——此时 $T$ 步之后 $\bw^{(T)}$ 近似于在某个 $\lambda(T)$ 下的岭回归解，且 $\lambda$ 随 $T$ 增大而减小。换到牛顿法，这个近似不成立：牛顿法每步都直接跳到当前二次模型的极小点，七步就到头（\S4.5），「早停」在这里等于没解完，而不是在做正则。

所以实践中的用法要分清：

\begin{center}
\small
\begin{tabular}{@{}lll@{}}
\toprule
\rowcolor{navy!7}
\tabhead 场景 & \tabhead 早停的作用 & \tabhead 说明\\
\midrule
梯度下降 / SGD & 近似 L2 正则 & 强度由 $T$ 与步长共同决定\\
牛顿法 / IRLS & 只是没算完 & 想正则就显式加 $\lambda$\\
分离数据 & 防止数值爆炸 & 治标，不治本\\
\bottomrule
\end{tabular}
\end{center}

因此「迭代次数」在工程上更该被当成\textbf{安全阀}：设一个上限，防止发散时无限循环；真正的正则应该写在目标函数里，而不是藏在循环次数里。

\begin{keybox}[本节结论]
\textbf{「无解」有三种成因，处置方式不同。}完全分离与拟完全分离是数据问题，加 L2、加 L1、或早停都能救；共线导致的 $\bH$ 奇异是设计问题，靠正则或降维；尺度悬殊是预处理问题，靠标准化。三者都表现为「不收敛」，但只有第一种是「真的没有解」。
\end{keybox}

% ============================================================
%  第 6 节
% ============================================================

\section{阈值、代价与校准}

\subsection{阈值的来源}

有了概率，还要把它翻译成类别。翻译规则就是 \S1.2 的式 \eqref{eq:bayesrule}，这里把它拆成两个尺度说清。

\begin{prop}[贝叶斯最优阈值]
沿用 \S1.2 的 $c_{ij}$（真类 $i$ 判成 $j$，$c_{00}=c_{11}=0$），则
\begin{equation}
  \text{判 }1\iff p(y=1\mid\bx)>\tau,\qquad \tau=\frac{c_{01}}{c_{01}+c_{10}} .
\end{equation}
等代价时 $\tau=1/2$。换到 $z$ 尺度，这条规则等价于
\begin{equation}
  \text{判 }1\iff z>z^{\ast},\qquad z^{\ast}=\mathrm{logit}(\tau)=\ln\frac{c_{01}}{c_{10}} .
  \label{eq:zstar}
\end{equation}
\end{prop}

\begin{proof}
$p>\tau\iff\mathrm{logit}(p)>\mathrm{logit}(\tau)$（$\mathrm{logit}$ 单调），而 $\mathrm{logit}\bigl(\frac{c_{01}}{c_{01}+c_{10}}\bigr)=\ln\frac{c_{01}}{c_{10}}$。
\end{proof}

两个尺度必须分开说，否则很容易写出自相矛盾的句子。要点是：

\begin{itemize}
  \item \textbf{$p$ 尺度的 $\tau$ 与先验无关}。因为 $p$ 本身就是后验，先验已经含在里面了。等代价时 $\tau=0.5$ 是一个\emph{结论}，不是约定。
  \item \textbf{先验的影响发生在 $z$ 尺度上，而且是平移截距}。设训练样本的类先验为 $(\phi_1,\phi_0)$，目标人群为 $(\phi_1',\phi_0')$，则
  \begin{equation}
    z'=z+\ln\frac{\phi_1'/\phi_0'}{\phi_1/\phi_0},
    \qquad\text{等价于}\qquad
    b'=b+\ln\frac{\phi_1'/\phi_0'}{\phi_1/\phi_0}.
    \label{eq:priorshift}
  \end{equation}
  直觉：目标人群里类 1 更多，同样的证据就该给更高的后验——截距上抬，边界整体往类 0 一侧挪。
\end{itemize}

式 \eqref{eq:priorshift} 与第二讲 \S2.7 的先验修正是同一件事，但\textbf{只借结构、不借符号}：第二讲的 $z$ 越大越像类 2，本讲的 $z$ 越大越像类 1，两边的平移量符号因此相反。写讲义时若直接抄公式，会得到一个「先验越大、越不像类 1」的荒谬结论。

\begin{warn}[阈值 $0.5$ 不是定理]
$0.5$ 只在「等代价 + 后验已校准」时成立。信用卡欺诈检测里漏判的代价是误判的上千倍，$\tau$ 应该远小于 $0.5$；而医学筛查里，若后续复查很贵，$\tau$ 又该往大调。把 $0.5$ 当默认值可以，当定理不行。
\end{warn}

\subsection{代价敏感与先验修正}

拿到不平衡数据，有两条路可以走，效果相近但代价不同。

\paragraph{阈值调整。}不动模型，只把判决规则里的 $\tau$ 换成式 \eqref{eq:bayesrule} 给出的值。优点是不用重训，改一行代码；缺点是这个 $\tau$ 依赖你对 $c_{01},c_{10}$ 的估计，而那往往比模型本身更难估。

\paragraph{样本加权。}给样本加权，等价于改变了训练分布，也等价于改变了「有效先验」。负类样本权重乘 $k$，与把正类样本复制 $k$ 倍是同一件事——两者都会把截距往正类一侧推。这在实现上就是给似然乘权重：
\begin{equation}
  L_w=-\frac{1}{\sum_i\omega_i}\sum_{i}\omega_i\,\ell_i .
\end{equation}
加权后的梯度与 Hessian 只是每个样本多乘一个 $\omega_i$，IRLS 的形式完全不变。

两者的区别值得记牢：\textbf{改阈值只改判决，改权重改模型}。如果下游还要用概率本身（\S6.3），就只能走第二条路——阈值调整不会让概率变得可信。

\begin{exm}[准确率 99\% 却毫无用处]
一千个样本，十个正类。模型把全部样本判为负类：准确率 $990/1000=99\%$，看起来漂亮；但召回率为 $0$，AUC 为 $0.5$（与随机猜没有区别）。这说明\textbf{准确率在不平衡数据上几乎不携带信息}——它把「全判多数类」这个平凡策略也给了高分。

换成 $\tau=0.05$ 之后，同一组概率可能给出 $80\%$ 的召回与 $15\%$ 的误报，准确率降到 $86\%$，但模型终于有用了。
\end{exm}

\subsection{校准与可靠性曲线}

\begin{dfn}[校准]
若对任意 $q\in(0,1)$，满足 $\hat p=q$ 的样本中真正属于类 1 的比例恰为 $q$，则称概率输出是\textbf{校准}的。
\end{dfn}

换句话说：$p=0.7$ 的那批样本里，应该真有约七成是正类。注意这与「排序对不对」是两件事——一个把概率全压在 $[0.4,0.6]$ 里的模型，排序可能完美（AUC 很高），但没人敢把 $0.55$ 解释成「五成半把握」。

校准的图形诊断是\textbf{可靠性曲线}：把预测概率分箱，横轴取箱内平均预测值，纵轴取箱内实际正类比例，理想情形是对角线。

\begin{figure}[htbp]
\centering
\begin{tikzpicture}
\begin{axis}[
  width=7.6cm,height=6.4cm,
  xlabel={预测概率 $\hat p$},ylabel={实际正类比例},
  xmin=0,xmax=1,ymin=0,ymax=1,
  axis lines=left,axis line style={steel!60,line width=0.5pt},
  tick style={steel!60},tick label style={font=\scriptsize,text=steel},
  label style={font=\scriptsize,text=steel},
  legend style={font=\scriptsize,draw=steel!30,fill=white,at={(0.03,0.97)},anchor=north west},
  grid=both,grid style={steel!12,line width=0.3pt},
]
\addplot[steel!55,line width=0.5pt,densely dashed] coordinates {(0,0)(1,1)};
\addlegendentry{理想（$y=x$）}
\addplot[gold,line width=0.9pt,mark=*,mark size=1.2pt] coordinates {
  (0.05,0.20)(0.15,0.30)(0.35,0.44)(0.65,0.56)(0.85,0.70)(0.95,0.80)};
\addlegendentry{未校准（过自信）}
\addplot[navy,line width=0.9pt,mark=o,mark size=1.2pt] coordinates {
  (0.05,0.06)(0.15,0.14)(0.35,0.34)(0.65,0.66)(0.85,0.84)(0.95,0.94)};
\addlegendentry{校准后}
\end{axis}
\end{tikzpicture}
\caption{可靠性曲线（示意数据）。未校准那条比对角线\textbf{平}：模型说 $0.95$，实际只有八成——这是过自信的典型形状。校准后重新贴回对角线。看图时只问一件事：曲线在哪些区间偏离对角线，偏的方向是「说大了」还是「说小了」。}
\label{fig:reliability}
\end{figure}

偏离对角线有\textbf{两个方向}，成因不同，别混成一句口号：

\begin{itemize}
  \item \textbf{过自信}（说 $0.95$ 实际 $0.8$）：典型来源是\textbf{近可分}与\textbf{稀疏单元}——某些系数被推向 $\pm\infty$，预测贴到 $0$ 或 $1$。
  \item \textbf{欠自信}（说 $0.6$ 实际 $0.9$）：典型来源是\textbf{强正则}——$\lambda$ 大时系数被压向零，概率被压向 $0.5$。
\end{itemize}

两者可以同时存在，谁占主导要看正则强度与样本量。这里要格外小心\textbf{「小样本必然过自信」这类口号}：有限样本偏差是 $O(1/n)$，方向随设计、稀疏程度甚至\emph{估计量的选法}而变——同一个 $p=0.9$、$n=50$ 的例子，剔除退化样本后的条件均值把 logit 推向远离 $0$（约 $2.29$，真值 $2.20$），而 Haldane 估计 $\mathrm{logit}\frac{x+0.5}{n+1}$ 几乎无偏（$2.20$）。换估计量，方向就变。所以只能说「机制不同」，不能说「必然朝哪边」。

修正方法常见两种：

\begin{itemize}
  \item \textbf{Platt 缩放}：拟合两个参数，$p\mapsto\sigma\bigl(A\,\mathrm{logit}(p)+B\bigr)$。$A>1$ 是锐化（治欠自信），$A<1$ 是压缩（治过自信），\textbf{两个方向都能修}；参数少，样本少吃得住。
  \item \textbf{isotonic 回归}：拟合一条单调不减的阶梯函数。它是\textbf{非参数}的，能修任意单调畸变，但自由度大，样本少时容易过拟合。
\end{itemize}

真正的区别在「参数化 vs 非参数」，而不是「各治一种畸变」。两者都要求一份独立的校准集——用训练集拟合校准参数会得到「校准得很好」的假象。

\subsection{评价指标的选择}

指标不是越多越好，而是\textbf{想比什么就选哪个}。四行表格把它们分开：

\begin{center}
\small
\begin{tabular}{@{}lp{4.2cm}p{4.0cm}p{3.4cm}@{}}
\toprule
\rowcolor{navy!7}
\tabhead 指标 & \tabhead 看什么 & \tabhead 依赖 & \tabhead 什么时候用\\
\midrule
准确率 & 判决对错的比例 & 阈值 & 类别均衡、代价相同\\
AUC & 排序（正类排在负类前的概率） & 只看排序，与校准无关 & 比较模型、类别不平衡\\
LogLoss & 概率的对数损失 & 排序 + 校准 & 需要概率输出时\\
Brier & 概率的平方误差 & 排序 + 校准 & 想要「概率误差」的直观量\\
\bottomrule
\end{tabular}
\end{center}

\begin{figure}[htbp]
\centering
\begin{tikzpicture}
\begin{axis}[
  width=8.4cm,height=7.0cm,
  xlabel={假正率 FPR},ylabel={真正率 TPR},
  xmin=0,xmax=1,ymin=0,ymax=1,
  axis lines=left,axis line style={steel!60,line width=0.5pt},
  tick style={steel!60},tick label style={font=\scriptsize,text=steel},
  label style={font=\small,text=steel},
  grid=both,grid style={steel!12,line width=0.3pt},
]
\addplot[steel!55,line width=0.5pt,densely dashed] coordinates {(0,0)(1,1)};
\addplot[navy,line width=1.0pt,mark=none] coordinates {
  (0,0)(0.03,0.35)(0.07,0.55)(0.12,0.68)(0.20,0.78)(0.32,0.86)(0.48,0.92)(0.70,0.97)(1,1)};
\addplot[gold,only marks,mark=*,mark size=2.2pt] coordinates {(0.12,0.68)(0.20,0.78)(0.48,0.92)};
\draw[steel!50,line width=0.3pt] (axis cs:0.48,0.92) -- (axis cs:0.53,0.80);
\draw[steel!50,line width=0.3pt] (axis cs:0.20,0.78) -- (axis cs:0.30,0.68);
\draw[steel!50,line width=0.3pt] (axis cs:0.12,0.68) -- (axis cs:0.22,0.56);
\node[font=\footnotesize,text=gold,anchor=west] at (axis cs:0.54,0.79) {$\tau=0.2$};
\node[font=\footnotesize,text=gold,anchor=west] at (axis cs:0.31,0.67) {$\tau=0.5$};
\node[font=\footnotesize,text=gold,anchor=west] at (axis cs:0.23,0.55) {$\tau=0.8$};
\end{axis}
\end{tikzpicture}
\caption{ROC 曲线与阈值的关系（示意数据）。曲线本身由全部可能的阈值连成，改 $\tau$ 只是在这条曲线上\textbf{移动一个点}：$\tau$ 小则往右上（多报、少漏），$\tau$ 大则往左下。AUC 是曲线下的面积，改阈值不改变它——这也说明 AUC 与「概率准不准」是两件事。}
\label{fig:roc}
\end{figure}

\begin{warn}[两个常见的误用]
\textbf{其一，把 $0.5$ 当定理}：在不平衡或代价不等的场合，$0.5$ 往往不是最优阈值（\S6.1）。

\textbf{其二，把 AUC 高当成概率准}：AUC 只依赖排序，把所有的 $p$ 换成 $p^{3}$（单调变换，排序不变），AUC 一点不变，但校准已经完全错了。要谈「$0.7$ 是不是七成把握」，只能看可靠性曲线与 LogLoss／Brier。
\end{warn}

\begin{keybox}[本节结论]
\textbf{输出的可信度要单独验证。}阈值由代价与先验决定，不由 $0.5$ 决定；校准由可靠性曲线检查，不由 AUC 代替；两者都不能靠训练集自我评估。一句话：模型给出 $0.7$，你要能回答「凭什么信这个 $0.7$」，否则它只是一个被压扁的得分。
\end{keybox}
% ============================================================
%  第 7 节
% ============================================================

\section{模型比较与选用}

\subsection{与感知机的差别}

逻辑回归与感知机共用同一个骨架 $\bw\tp\bx$，判决规则也都是「看符号」。差别全在训练目标上：

\begin{center}
\small
\begin{tabular}{@{}lp{5.4cm}p{5.4cm}@{}}
\toprule
\rowcolor{navy!7}
\tabhead & \tabhead 感知机 & \tabhead 逻辑回归\\
\midrule
损失 & 合页式：$\max(0,-s_i z_i)$ & 交叉熵：$\ln(1+e^{-s_iz_i})$\\
光滑性 & 分段线性，不可微 & 无穷次可微\\
解 & 非凸；可分时解不唯一 & 凸；非退化时唯一\\
输出 & 符号 $\pm1$ & 概率 $p\in(0,1)$\\
\bottomrule
\end{tabular}
\end{center}

两行值得展开。其一，\textbf{合页损失只惩罚「错得不够远」的样本}：一旦 $s_iz_i>1$，该样本的梯度就是零，模型不再看它；交叉熵则在所有样本上都有非零梯度（$p_i-y_i\ne0$），只是随距离指数衰减。这解释了逻辑回归对「已经分对的边缘样本」仍会微调，而感知机不会。

其二，\textbf{输出语义不同}。感知机给的是「哪边」，逻辑回归给的是「多大概率在哪边」。后者能接 \S6 的全部机制（代价、阈值、校准），前者只能接「错几个」。

\subsection{与 LDA 的取舍}

第二讲 \S4.5 有一张 LDA 与逻辑回归的对照表，第三讲 \S8.3 又补了一张 GDA 与逻辑回归的（把样本效率拆成「样本少」「样本多」两行，另加「假设不成立」一栏），这里不重复。本节只补那两张表里没说的两件事。

\paragraph{收敛速率与渐近误差。}常被引用的定量结论来自 Ng \& Jordan (2001)。读这篇时最容易错的两个地方，先标出来：

\begin{itemize}
  \item 文中的量是\textbf{参数个数} $n$（在本讲即 $D+1$），\textbf{不是样本数} $N$；
  \item 「渐近有效」说的是\emph{参数估计}达到 Cramér--Rao 下界，\emph{不是}说分类误差达到下界。
\end{itemize}

在这个前提下，结论可以分成两句：

\begin{itemize}
  \item \textbf{样本复杂度}：需要多少样本才能「接近自己的极限」，生成式随 $n$ 只呈\textbf{对数}增长，判别式呈\textbf{线性}增长。所以小样本时 GDA 更省。
  \item \textbf{渐近误差}：若高斯假设成立，GDA 对参数估计是渐近有效的，逻辑回归一致但渐近方差更大；若假设不成立，则反过来——GDA 的渐近误差更高，逻辑回归占优。
\end{itemize}

两句话合起来，就是那条著名的「交叉」：小样本时生成式领先，样本变多后判别式追上并反超。所以「逻辑回归需要 $\log N$ 倍的样本」是误读——正确的说法是「两者的样本量需求分别随参数个数线性增长与对数增长」。

\paragraph{诊断式选型。}不要背口号，看两个量：

\begin{enumerate}[label=\textcolor{navy}{\bfseries\arabic*.},leftmargin=2.0em]
  \item \textbf{协方差是否齐性}：逐类算 $\abs{\bS_c}$ 与 $\tr(\bS_c)$，看是否同量级（第二讲 \S6.5 第 1 条）。差异大说明各类形状不同，GDA 的前提破了。
  \item \textbf{$N/D$ 够不够}：共享协方差要估 $D(D+1)/2$ 个参数。$N$ 只是 $D$ 的两三倍时，估计的方差会很大，此时哪怕高斯假设成立，LDA 也未必赢。
\end{enumerate}

\textbf{正态性不作门槛。}第二讲 \S6.5 的注意框已经说明：即使数据不正态，$\hat\bw\propto\bS_w^{-1}(\bar\bx_2-\bar\bx_1)$ 仍是一个合理方向，正态假设只在「它是否还是最优贝叶斯分类器」这一层起作用。若确实想查，用偏度、QQ 图或 Mardia 检验，但它决定的是\emph{效率}，不是\emph{可用性}。

\begin{keybox}[怎么选]
\textbf{小样本 + 协方差齐性 OK}：先试 LDA，它省样本。\\
\textbf{样本充足 / 协方差明显不齐 / 只关心 $p(y\mid\bx)$}：用逻辑回归，它更抗假设错误，也直接给概率。\\
\textbf{两者都跑一遍}：交叉验证比一比，这比任何先验判断都可靠。
\end{keybox}

\subsection{失效场景清单}

按踩坑频率排：

\begin{enumerate}[label=\textcolor{navy}{\bfseries\arabic*.},leftmargin=2.0em,itemsep=0.45em]
  \item \textbf{完全分离／拟完全分离}：系数发散，\texttt{max\_iter} 报收敛失败。处置：加 L2（\S5.3）、加 L1、或换 Firth 惩罚似然。
  \item \textbf{强共线}：$\bH$ 近奇异，牛顿步乱跳。处置：去掉冗余特征、加 $\epsilon\bI$、改 L-BFGS。
  \item \textbf{特征尺度悬殊}：$\bH$ 条件数爆炸。处置：标准化（且标准化参数只能来自训练集）。
  \item \textbf{极端不平衡 + 未校准}：准确率虚高，概率不可信。处置：调阈值或加权（\S6.2），并做校准（\S6.3）。
  \item \textbf{$p$ 集中在 $0/1$}：梯度 $p_i-y_i$ 趋于零，训练几乎不动。处置：检查是否已分离，或改用带间隔的损失。
  \item \textbf{把它当非线性模型用}：逻辑回归只能给出\textbf{线性}边界（在原始特征空间里）。要非线性，先做基函数展开（第一讲 \S6.1），再把展开后的特征喂进来。
\end{enumerate}

最后一条尤其值得强调：模型的「非线性能力」不在链接函数里，而在特征里。sigmoid 只负责把实数压成概率，它不负责把弯曲的边界拉直。

% ============================================================
%  第 8 节
% ============================================================

\section{小结}

\subsection{五个要点}

\begin{enumerate}[label=\textcolor{navy}{\bfseries\arabic*.},leftmargin=2.3em,itemsep=0.55em]
  \item \textbf{逻辑回归 = 线性打分 + sigmoid，而 sigmoid 是反解出来的。}假设只有一条：对数几率线性。它可能来自高斯假设，也可以独立成立。
  \item \textbf{损失就是平均交叉熵，也就是负对数似然。}三者与 KL 最小等价；梯度是「残差 × 特征」，与最小二乘同构。
  \item \textbf{严格凸，所以解唯一。}牛顿法重排之后就是 IRLS——每一步都是一次加权最小二乘，第一讲的机器原样复用。
  \item \textbf{分离时无解，正则化把解拉回来。}完全分离与拟完全分离都会让 $\norm{\bw}\to\infty$；L2 等价于高斯先验的 MAP，L1 兼做特征选择。
  \item \textbf{概率要校准才算「把握」。}阈值由代价与先验决定，校准由可靠性曲线验证，而 AUC 两件事都不管。
\end{enumerate}

\subsection{知识结构}

\begin{figure}[htbp]
\centering
\begingroup
\providecommand{\ttl}{\bfseries\sffamily\color{navy}}
\providecommand{\fml}{\footnotesize\color{navy!88}}
\scalebox{1.0}{%
\begin{tikzpicture}[
  box/.style={draw=navy!55,fill=navy!3,minimum width=2.62cm,minimum height=1.92cm,
    align=center,inner sep=2pt,sharp corners,line width=0.5pt},
  ar/.style={->,navy!60,line width=0.6pt},
]
\node[box] (b1) at (1.45,0) {{\ttl 模型}\\[2pt]{\fml $z=\bw\tp\bx$}\\[0pt]{\fml $p=\sigma(z)$}};
\node[box] (b2) at (4.45,0) {{\ttl 准则}\\[2pt]{\fml $L=-\frac1N\sum_i\ell_i$}\\[0pt]{\fml 交叉熵 = KL}};
\node[box] (b3) at (7.45,0) {{\ttl 求解}\\[2pt]{\fml 牛顿 = IRLS}\\[0pt]{\fml $\bH$ 与 $\bW$}};
\node[box] (b4) at (10.45,0) {{\ttl 存在性}\\[2pt]{\fml 分离 $\Rightarrow$ 无解}\\[0pt]{\fml L2 = MAP}};
\node[box] (b5) at (13.45,0) {{\ttl 使用}\\[2pt]{\fml $\tau=\frac{c_{01}}{c_{01}+c_{10}}$}\\[0pt]{\fml 校准 + 指标}};
\draw[ar] (b1) -- (b2);
\draw[ar] (b2) -- (b3);
\draw[ar] (b3) -- (b4);
\draw[ar] (b4) -- (b5);
\draw[<->,navy!45,line width=0.55pt,densely dashed] (b1.north) to[bend left=24] (b3.north);
\draw[<->,navy!45,line width=0.55pt,densely dashed] (b2.south) to[bend right=24] (b4.south);
\node[draw=navy!25,fill=white,align=center,font=\footnotesize,inner sep=5pt]
  at (7.45,-2.34)
  {判别式：只管 $p(y\mid\bx)$ \quad $\cdot$ \quad 生成式：先管 $p(\bx\mid y)$ \quad $\cdot$ \quad 差别在假设，不在骨架};
\end{tikzpicture}}
\endgroup
\caption{本讲的知识地图。横向是一条流水线：模型、准则、算法、存在性、使用。两条虚线是本讲反复出现的回路——求解的困难会反过来迫使你改准则（分离 $\to$ 加正则），而使用中的偏差又会迫使你回头检查模型（校准）。}
\label{fig:map}
\end{figure}

\subsection{后续内容}

这一讲停在「线性边界 + 校准过的概率」上。往下走有三条路：

\begin{itemize}
  \item \textbf{多分类}。把 sigmoid 换成 softmax，把伯努利似然换成多项似然，就是多项逻辑回归。$K$ 类交叉熵、参数的可辨识性、一对一与一对多策略，都是独立的话题，本讲刻意没有展开。
  \item \textbf{非线性边界}。基函数展开（第一讲 \S6.1）最省事，核方法与神经网络更彻底。注意：换了特征之后，模型仍然是「线性打分 + sigmoid」，只是打分发生在展开后的空间里。
  \item \textbf{贝叶斯视角}。给 $\bw$ 一个先验，用 Laplace 近似或 MCMC 求后验，就得到「参数的不确定性」——这比一个点估计更适合小样本与在线决策。
\end{itemize}

% ============================================================
%  习题
% ============================================================

\clearpage
\section*{习题}
\addcontentsline{toc}{section}{习题}
\markboth{习题}{}
\setlength{\parskip}{0.62em}

\subsection*{A 组 · 基础}

\noindent\textbf{A1.} 用 \S4.5 的解 $\hat{\bw}=(-1.0307,\,-2.6288,\,8.2214)$，手算第 4 点 $(2.6,2.1)$ 的 $z$、$p$ 与 log-odds，并解释为什么 $p=0.5053$ 会被判为类 1。

\noindent\textbf{A2.} 取 $p=0.9$，算出对应的 $z$；再用 $z'=-z$ 反解出 $p'$，验证 $\sigma(-z)=1-\sigma(z)$。用同一个 $z$ 算 $\sigma'(z)$，说明为什么它最大不超过 $1/4$。

\noindent\textbf{A3.} 写出 $y\in\{-1,1\}$ 下的对数似然 $-\ln\bigl(1+e^{-y z}\bigr)$，并说明为什么把 $\pm1$ 直接代进 $p^{y}(1-p)^{1-y}$ 会失效。

\noindent\textbf{A4.} 交叉熵与 KL 散度相差的那一项是什么？在本讲的例题上把它算出来（提示：两类样本各三个，经验分布的熵只与频率有关）。

\subsection*{B 组 · 提高}

\noindent\textbf{B1.} 证明 $\bH=\frac1N\bX\tp\bW\bX\succeq0$，并给出「正定」的充要条件。举一个 $\bH$ 半正定但不满秩的例子。

\noindent\textbf{B2.} 证明 IRLS 的一步（式 \eqref{eq:irls}）等价于以 $\bW$ 为权的加权最小二乘。提示：从 $\bX\tp\bW\bX\,\bDelta=\bX\tp(p-y)$ 出发，把 $\bDelta$ 移项。

\noindent\textbf{B3.} 构造一个一维的完全分离例子，说明 $\hat w$ 不存在；再构造一个拟完全分离的例子，并验证：若允许 $\bw=\bzero$，定义 (ii) 对\emph{任何}数据集都成立，所以「$\bw\ne\bzero$」这个条件是必需的。

\noindent\textbf{B4.} 用 \S4.5 的六点数据验证：边界 $x+y=3$（判 1 当 $x+y<3$）只错第 3 点，而交叉熵最优解错两个点。算出两条边界的错分个数与交叉熵，解释为什么「交叉熵最小」不等于「错分最少」。

\subsection*{C 组 · 动手}

\noindent\textbf{C1.} 不调用现成库，实现 \S4.3 的 IRLS，在 \S4.5 的六点上与 \texttt{sklearn} 的实现对齐到 $10^{-8}$（它默认带 L2 惩罚，记得显式关掉）。

\noindent\textbf{C2.} 在删掉第 3 点的可分数据上，把 \texttt{max\_iter} 从 $10$ 加到 $500$，画出 $\norm{\bw^{(t)}}$ 随 $t$ 的变化，判断它是线性增长还是指数增长；再把 $p$ 打印出来，找出第一次出现 \texttt{nan} 的迭代步。

\noindent\textbf{C3.} 用 Platt 缩放（式中的 $\sigma(A\,\mathrm{logit}(p)+B)$）对一组过自信的概率做校准，报告 $A$ 落在 $1$ 的哪一侧，并画出校准前后的可靠性曲线。

\subsection*{提示}

\begin{itemize}[itemsep=0.2em]
  \item \textbf{A1}：$z=\bw\tp\bx$ 里别忘了截距项；$z>0$ 就判 1。
  \item \textbf{A2}：$\sigma'(z)=\sigma(z)(1-\sigma(z))$ 是开口向下的抛物线，顶点在 $p=1/2$。
  \item \textbf{A4}：经验分布是 $\{1/2,1/2\}$（两类各半），其熵为 $\ln2\approx0.693$；用它减去 $L=0.3696$。
  \item \textbf{B1}：$\bv\tp\bH\bv=\frac1N\sum p_i(1-p_i)(\bv\tp\bx_i)^2$，正定要求所有权重为正且 $\bX$ 列满秩。
  \item \textbf{B2}：加权最小二乘的目标是 $\sum \omega_i(\text{残差}_i)^2$，其正规方程与 \eqref{eq:newtonstep} 同形。
  \item \textbf{B3}：拟完全分离的例子里，让某个特征在正类全为 $1$、负类全为 $0$，其余特征相同。
  \item \textbf{C2}：$\bW$ 的元素会下溢到 $0$，此时 $\bH$ 奇异，先看这一步之前发生了什么。
\end{itemize}

% ============================================================
%  附录
% ============================================================

\appendix
\renewcommand{\thesection}{\Alph{section}}

\section{公式速查}

\subsection*{A.1 三个函数}
\begin{align}
  \mathrm{logit}(p)&=\ln\frac{p}{1-p}, &
  \sigma(z)&=\frac{1}{1+e^{-z}}, &
  \sigma'(z)&=\sigma(z)\bigl(1-\sigma(z)\bigr),\\
  \sigma(-z)&=1-\sigma(z), &
  -\ln\sigma(z)&=\softplus(-z), &
  -\ln\bigl(1-\sigma(z)\bigr)&=\softplus(z).
\end{align}
其中
\begin{equation}
  \softplus(z)=\ln(1+e^{z})=\max(z,0)+\ln\bigl(1+e^{-\abs{z}}\bigr).
\end{equation}

\subsection*{A.2 梯度、Hessian 与 IRLS}
\begin{align}
  \bg&=\frac1N\sum_{i=1}^{N}(p_i-y_i)\,\bx_i, &
  \bH&=\frac1N\bX\tp\bW\bX,\quad \bW=\diag\bigl(p_i(1-p_i)\bigr),\\
  \bw&\leftarrow\bw-\bH^{-1}\bg, &
  \bw&\leftarrow\bigl(\bX\tp\bW\bX\bigr)^{-1}\bX\tp\bW\bz_{\mathrm w},\\
  \bz_{\mathrm w}&=\bX\bw+\bW^{-1}(y-p), &
  \bH_{\lambda}&=\bH+\lambda\begin{pmatrix}\bI_D&\\&\mathbf{0}\end{pmatrix}.
\end{align}

\subsection*{A.3 阈值与代价}
\begin{align}
  \tau&=\frac{c_{01}}{c_{01}+c_{10}}, &
  z^{\ast}&=\ln\frac{c_{01}}{c_{10}}, &
  b'&=b+\ln\frac{\phi_1'/\phi_0'}{\phi_1/\phi_0}.
\end{align}

\subsection*{A.4 四个指标}
\begin{itemize}
  \item 准确率：$\dfrac{1}{N}\sum_i\mathbb{I}\bigl(\hat y_i=y_i\bigr)$，依赖阈值。
  \item AUC：随机取一正一负，正类得分更高的概率；只依赖排序。
  \item LogLoss：$\dfrac1N\sum_i\bigl[y_i\ln p_i+(1-y_i)\ln(1-p_i)\bigr]$ 的相反数，即本讲的 $L$。
  \item Brier：$\dfrac1N\sum_i(p_i-y_i)^2$，概率的平方误差。
\end{itemize}

\subsection*{A.5 记号对照}
\begin{center}
\small
\begin{tabular}{@{}lll@{}}
\toprule
\rowcolor{navy!7}
\tabhead 本讲 & \tabhead 常见别名 & \tabhead 说明\\
\midrule
$z=\bw\tp\bx$ & 线性预测子、logit & 笔记里的 $a$ 即此 $z$\\
$L$ & 交叉熵、负对数似然 & 平均尺度，越小越好\\
$\tau$ & 判决阈值 & $p$ 尺度；$z$ 尺度记 $z^{\ast}$\\
$\bDelta$ & 牛顿增量 & 不用 $\delta$：$\delta_c$ 是判别函数\\
$\bz_{\mathrm w}$ & 工作响应 & $N$ 维向量，与 $z_i$ 不同\\
$\bW$ & 权重矩阵（$N\times N$） & $=\diag(\Var y_i)$；第二讲的 $\bW$ 是投影矩阵，别混\\
$\lambda$ & 正则强度 & 与学习率 $\eta$ 分开，不再共用一个符号\\
\bottomrule
\end{tabular}
\end{center}

\section{延伸阅读}

\begin{enumerate}[label={[\arabic*]},leftmargin=2.6em,itemsep=0.7em]
  \item T. Hastie, R. Tibshirani, J. Friedman. \emph{The Elements of Statistical Learning}, 2nd ed. Springer, 2009.\\[1pt]
  {\footnotesize\color{steel}\S4.4 节：逻辑回归的完整处理，含 IRLS、分离现象与 $\ell_1$ 正则。本讲 \S4 的骨架与它一致。}

  \item A. Y. Ng, M. I. Jordan. On discriminative vs.\ generative classifiers: a comparison of logistic regression and naive Bayes. \emph{NIPS}, 2001.\\[1pt]
  {\footnotesize\color{steel}「样本复杂度随参数个数对数／线性增长」的出处。注意原文对照的是朴素贝叶斯，GDA 是同一框架下的连续版本。}

  \item A. Albert, J. A. Anderson. On the existence of maximum likelihood estimates in logistic regression models. \emph{Biometrika}, 71(1):1--10, 1984.\\[1pt]
  {\footnotesize\color{steel}完全分离、拟完全分离与 MLE 存在性的完整分类。本讲 \S5.1 的判据来自此文。}

  \item D. Firth. Bias reduction of maximum likelihood estimates. \emph{Biometrika}, 80(1):27--38, 1993.\\[1pt]
  {\footnotesize\color{steel}用 Jeffreys 先验的一阶修正，既减小偏差，又能在分离时给出有限解。是 \S5.3 之外的另一条补救路线。}

  \item J. Platt. Probabilistic outputs for support vector machines and comparisons to regularized likelihood methods. \emph{Advances in Large Margin Classifiers}, 1999.\\[1pt]
  {\footnotesize\color{steel}Platt 缩放的原始文献。两参数、样本少也稳，是 \S6.3 的首选基线。}

  \item S. Boyd, L. Vandenberghe. \emph{Convex Optimization}. Cambridge University Press, 2004.\\[1pt]
  {\footnotesize\color{steel}\S9.5 与 \S10.2：牛顿法、收敛阶与自协和性。想弄清「二次收敛」的确切含义，看这里。}

  \item \texttt{scikit-learn} 文档：\texttt{LogisticRegression}。\\[1pt]
  {\footnotesize\color{steel}参数 \texttt{penalty}、\texttt{C}（$=1/\lambda$）、\texttt{class\_weight} 与 \texttt{max\_iter} 的确切含义，以及它内部用的求解器。}
\end{enumerate}

\vspace{1.0em}
\begin{flushright}
{\color{navy!25}\rule{3.2cm}{0.5pt}}\\[5pt]
{\kaishu\footnotesize\color{steel}CyArk}
\end{flushright}

\end{document}
